\documentclass[11pt]{article}
\usepackage[a4paper,margin=1.05in]{geometry}
\usepackage{amsmath,amssymb,amsthm,mathtools}
\usepackage{bm,microtype,enumitem,booktabs,array,xcolor,mathrsfs}
\usepackage[hidelinks]{hyperref}
\usepackage{aliascnt}
\usepackage[nameinlink,capitalize]{cleveref}
\usepackage[backend=biber,style=alphabetic,maxnames=99]{biblatex}
\addbibresource{references_draft16.bib}

\newtheorem{theorem}{Theorem}[section]

\newaliascnt{proposition}{theorem}
\newtheorem{proposition}[proposition]{Proposition}
\aliascntresetthe{proposition}

\newaliascnt{lemma}{theorem}
\newtheorem{lemma}[lemma]{Lemma}
\aliascntresetthe{lemma}

\newaliascnt{corollary}{theorem}
\newtheorem{corollary}[corollary]{Corollary}
\aliascntresetthe{corollary}

\newaliascnt{claim}{theorem}
\newtheorem{claim}[claim]{Claim}
\aliascntresetthe{claim}

\newaliascnt{hypothesis}{theorem}
\newtheorem{hypothesis}[hypothesis]{Hypothesis}
\aliascntresetthe{hypothesis}

\theoremstyle{definition}
\newaliascnt{definition}{theorem}
\newtheorem{definition}[definition]{Definition}
\aliascntresetthe{definition}

\newaliascnt{remark}{theorem}
\newtheorem{remark}[remark]{Remark}
\aliascntresetthe{remark}

\newaliascnt{status}{theorem}
\newtheorem{status}[status]{Status note}
\aliascntresetthe{status}

\crefname{theorem}{Theorem}{Theorems}
\Crefname{theorem}{Theorem}{Theorems}
\crefname{proposition}{Proposition}{Propositions}
\Crefname{proposition}{Proposition}{Propositions}
\crefname{lemma}{Lemma}{Lemmas}
\Crefname{lemma}{Lemma}{Lemmas}
\crefname{corollary}{Corollary}{Corollaries}
\Crefname{corollary}{Corollary}{Corollaries}
\crefname{claim}{Claim}{Claims}
\Crefname{claim}{Claim}{Claims}
\crefname{hypothesis}{Hypothesis}{Hypotheses}
\Crefname{hypothesis}{Hypothesis}{Hypotheses}
\crefname{definition}{Definition}{Definitions}
\Crefname{definition}{Definition}{Definitions}
\crefname{remark}{Remark}{Remarks}
\Crefname{remark}{Remark}{Remarks}
\crefname{status}{Status note}{Status notes}
\Crefname{status}{Status note}{Status notes}

\newcommand{\e}{\mathrm e}
\newcommand{\1}{\mathbf 1}
\newcommand{\E}{\mathcal E}
\newcommand{\R}{\mathbb R}
\newcommand{\C}{\mathbb C}
\newcommand{\Z}{\mathbb Z}
\newcommand{\norm}[1]{\left\lVert #1\right\rVert}
\newcommand{\soft}{X^{o(1)}}

\title{\textbf{A Source-Explicit Conditional Common-Weight All-$m$ Type-II Estimate\\
and the Complementary Signed High-Order Frontier}\\
\large Late 27 August snapshot: double-GCD B\'ezout reduction, rank-one M\"obius source, and high-profile inverse frontier}
\author{Research draft}
\date{August 27, 2026}

\begin{document}
\maketitle

\begin{abstract}
This note records an audited public-facing state of a parity-sensitive Type-II
programme.  Part~I studies a canonical common-weight direct operator.  A row
is indexed by a prime $r\asymp R$, a shift $k$, and $m\asymp M$, with
$m'=m+kr\asymp M$.  Under five explicit source/analytic pins---one-sided
Poisson--Bruhat transport and an absolute source ceiling, fixed-state
exclusion, the outer-root energy pin, projective source routing, and the
literal recombination estimate---band-limited prime participation yields the
required $R^{-1/4+o(1)}$ contraction.  The finite participation inequality,
projective ratio-class algebra, exponent ledger, and normalization compilers
are separately machine-checked; the displayed source/analytic pins are not
silently promoted into Lean theorems.

Part~II records the complementary Gate-1B frontier.  Returning to the
native order-five shell and opening the first $TT^\ast$ covariance gives an
exact sector law: the same-modulus branch is a character Gram, the coprime
cross-modulus branch is a CRT Kloosterman family, and the shared-divisor branch
belongs to the existing GCD router.  The proper smooth same-$q$ cross-character
child and the coprime cross-$q$ no-wrap operator remain useful positive
analytic children.  The later high-conductor audits recovered the primitive
projector algebra, proved the primitive-Kloosterman inversion, and retired the
naive square-complement Pascadi closure route.  Late update~VI identified the
source-faithful native $4$M$\mid5$D signed MAM45 parent, and Late update~VII
reduced its clean $|J|=5$ face to a coherent signed projector.  Late
update~VIII is now controlling.  In the entropy-gap fully-crossed child, a
second pair of GCD reductions gives a B\'ezout line; large common-GCD tails are
logarithmically disposable and the primitive $g=s=1$ child carries the entire
natural $QV$ scale.  Rigidity alone is therefore certified nonclosing.

The primitive residual has the exact line form
\[
 \ell A_0-u w_0=2,\qquad
 A_t=A_0+ut,\qquad w_t=w_0+\ell t,
\]
and the clean semiprime coefficient simplifies to
\[
 \beta_{D,P}(n)=-\mu(n)L_{D,P}(n),\qquad
 L_{D,P}(n)=\sum_{\substack{p\mid n\\p\sim P,\ n/p\sim D}}\log p.
\]
Thus the parity carrier in the hard primitive child is literally
$\mu(A_t)\mu(A_{t+k})$, with only $O(1)$ large-prime selector multiplicity.
The first exact analytic residual is the rank-one ordinary correlation
\[
 \boxed{\texttt{RANKONE-BEZOUT-AFFINE-ELLIOTT5D45}},
\]
whose target is
\[
 \sum_{u,\ell,t,k\ne0}
 \mu(A_t)\mu(A_{t+k})
 L_{D,P}(A_t)L_{D,P}(A_{t+k})
 b_5(w_t)\overline{b_5(w_{t+k})}\,\mathcal W
 \ll_A QV(\log X)^{-A}.
\]
The shift family is rank one: $(r,s)=(uk,\ell k)$, so the actual shift family
has size $URH=X$, whereas the independent shift rectangle has size
$QV=XH$ with $H\ge X^{5/18}$.  Consequently published complete-shift
averaged-Chowla results do not splice directly without a fixed-power tax.
The low-conductor/low-frequency pretentious profile is internally eliminated
by the centred five-defect source and small-conductor prime/M\"obius
uniformity.  What survives is a high-conductor, high-frequency, or distributed
line-spectrum inverse problem.  Source-preserving MRK5--MOYK5 manufacture,
Wright's partially fixed denominator theorem, and Blomer--Pascadi's all-moduli
Kloosterman theorem have all been literally auditioned; none closes the
physical range.  The best current bound for the primitive B\'ezout residual is
$QV(\log X)^{C_0}$ against the required $QV(\log X)^{-A}$: there is no
remaining power deficit, but the arbitrary-log cancellation theorem is still
open.  Comparison/main-term matching, lower-defect packets, the $r>1$
square-character family, the transition strip, QK56 parent/exhaustiveness,
and finite source reassembly remain separate obligations.  Gate~1B is not
closed.

Part~III records a refined Ford--Maynard landing strip.  Ford--Maynard
Theorem~8.3 gives a published five-collar description of the exceptional
$H_2$ geometry, and Definition~5.1 identifies the coordinates of $\mathbf v(n)$
with logarithms of the individual prime factors.  This permits a more
source-specific treatment of the rough $N_2$ sector: squareful terms are
negligible, while the squarefree shifted-prime lane reduces, at the research
level, to a standard uniform upper sieve for two linear forms plus the published
rough-number geometry.  A proof audit of Proposition~7.22 shows that the
published universal Type-II hypothesis is invoked only after the proof has
generated the grouped equation~(7.23), so controlling every \emph{actual
generated} instance of that structured sum is a proof-specific sufficient
landing point.  Internal Ford-side archaeology further indicates finite-depth
generated one-variable factors and finite nuclear/weight-dependence compilers;
however, the literal source-exact weighted high-complexity packet dictionary
from Gate outputs to the Ford-generated factors remains missing.  We call this
source interface \texttt{SOURCE-ADAPTER45}.  These are programme reductions,
not weaker theorems published by Ford--Maynard.

No claim of Gate-1B closure, Full Type-II reassembly, a Hardy--Littlewood
prime-pair asymptotic, or twin-prime infinitude is made.
\end{abstract}

\tableofcontents
\section{Introduction}
A recurring obstruction in parity-breaking sieve arguments is that strong
fixed-slice estimates need not survive coherent averaging over the true source
family. The correct object is therefore the source-restricted family operator,
not a collection of unrelated local norms.

This note isolates a canonical direct Type-II operator and derives a conditional
all-$m$ mean-square bound under explicit source and analytic pins. Four structural ingredients drive the proof:
\begin{enumerate}[label=(\roman*)]
  \item exact source and determinant-$2$ row geometry;
  \item source-preserving Fourier/Poisson transport to a smooth moving-$r$
        family;
  \item band-limited prime participation, converting smoothness in $r$ into a
        moving-family contraction;
  \item an exact projective crossed-convolution identity for the zero-curvature
        branch.
\end{enumerate}

The downstream motivation is the Ford--Maynard prime-producing sieve
framework, where Type I and Type II information for an actual non-negative
sequence is used simultaneously \cite{FordMaynard2024}. Their work also shows
that a substantial Type II range is genuinely necessary for non-trivial lower
bounds. We do not perform the full downstream reassembly here.

\subsection{External analytic input}
The participation step needs a prime number theorem in intervals of length
$R^{3/4}$.  Huxley's classical short-interval prime number theorem gives,
uniformly for $x\asymp R$,
\[
   \pi(x+H)-\pi(x)=(1+o(1))\frac{H}{\log x}
   \qquad(H\ge x^{7/12+\varepsilon}),
\]
for every fixed $\varepsilon>0$ \cite{Huxley1972}.  Since $3/4>7/12$, the
plateaux used below are longer than the range required by Huxley's theorem,
so no modern strengthening is needed.  We also use the classical Bernstein
inequality
\[
   \|P'\|_\infty\le N\|P\|_\infty
\]
for trigonometric polynomials of degree at most $N$ \cite{Zygmund2002}.

\subsection{Scope}
The theorem proved below concerns the canonical \emph{common-weight} direct
source.  In the authoritative source, the physical smoothing datum is one
weight $W_D$; row dependence that appears later is derived deterministically
from this common datum through Fourier, reciprocity and Poisson transforms.
An abstract interface allowing arbitrary independent edge-dependent weights
is strictly more general and belongs to an upstream source-adaptation problem.
Likewise, a global packet census and Full Type II reassembly are not part of
the theorem proved here.  This distinction is also machine-checked in the
formal development: the Gate-1A direct closure certificate contains no global
packet-dictionary or edge-dependent-D2 field.

\section{Parameters and notation}
Let $X\to\infty$. Set
\[
    M=X^{1/3},\qquad R=X^a,\qquad L=X^b,
\]
and
\[
    H=X^{a+2b-2/3},\qquad
    K=\frac{M}{R},\qquad
    U=\frac{L}{H},\qquad
    D=\frac{L^2}{H}.
\]
We repeatedly use
\begin{equation}
    RK=M,\qquad DH=L^2,\qquad M^2H=RL^2.
    \label{eq:scale-identities}
\end{equation}
The three controlling vertices are
\[
   V_1=\left(\frac{5}{18},\frac13\right),\qquad
   V_2=\left(\frac{5}{18},\frac{25}{72}\right),\qquad
   V_3=\left(\frac{7}{24},\frac13\right).
\]

For a prime $\ell$, define
\begin{equation}
   \rho_\ell(n):=\1_{\ell\mid n}-\frac1\ell.
   \label{eq:rho-def}
\end{equation}
Let $b_p,d_q$ be $1$-bounded coefficient sequences on primes $p,q\asymp L$.
All harmless divisor and logarithmic losses are absorbed in $X^{o(1)}$.

We use $e(z)=e^{2\pi iz}$ and $e_c(z)=e(z/c)$.  For a Schwartz function
$f$ we use the Fourier convention
\[
   \widehat f(\xi)=\int_{\mathbb R} f(x)e(-x\xi)\,dx.
\]
Every inverse symbol in a modular phase is understood only on the explicitly
routed unit/full-conductor sector.

\section{The canonical physical source}
\subsection{Rows and affine forms}
A canonical row is $e=(r,k,m)$ with
\[
   r\asymp R \text{ prime},\qquad 1\le |k|\le K,\qquad
   m\asymp M,\qquad m':=m+kr\asymp M,
\]
together with the generic unit/full-conductor conditions needed by the
reciprocal phases. Let $\E^\sharp$ be this row family.

Choose $w_0$ modulo $r$ such that
\[
   mw_0\equiv -2\pmod r.
\]
Since $m'=m+kr\equiv m\pmod r$, the same choice satisfies
$m'w_0\equiv-2\pmod r$.  Hence
\[
   \alpha:=\frac{mw_0+2}{r},\qquad
   \beta:=\frac{m'w_0+2}{r}
\]
are integers.  Define
\[
   A_e(t):=\alpha+mt,\qquad
   B_e(t):=\beta+m't.
\]

\begin{lemma}[Determinant-$2$ row identity]
For every canonical row and every integer $t$,
\[
   m'A_e(t)-mB_e(t)=2k.
\]
\end{lemma}
\begin{proof}
The $t$-terms cancel and
\[
  m'\alpha-m\beta
  =\frac{m'(mw_0+2)-m(m'w_0+2)}{r}
  =\frac{2(m'-m)}r=2k.
\]
\end{proof}

\subsection{Common physical smoothing}
Fix $W\in C_c^\infty(\R)$ and write $W_D(t)=W(t/D)$. The key point is that
this weight is common to all rows. Define
\[
   P_e(t)=\sum_{p\asymp L}b_p\rho_p(A_e(t)),\qquad
   Q_e(t)=\sum_{q\asymp L}d_q\rho_q(B_e(t)),
\]
and
\begin{equation}
   C_e(b,d):=\sum_{t\in\Z}W_D(t)P_e(t)Q_e(t).
   \label{eq:Ce}
\end{equation}

\begin{remark}
An arbitrary family $W_{D,e}$ is not part of the canonical theorem. Such
weights arise only in a more general source-adaptation interface.
\end{remark}

\subsection{Source provenance and derived row dependence}
The common-weight assertion is stronger than a notational convention.  After
Fourier transformation the source coefficient factors as
\begin{equation}
  \omega_{m,p,q}(h)
  = \underbrace{\frac{H}{pq}
      \e\!\left(-\frac{h\alpha}{pqm}\right)}_{\text{derived row-dependent scalar}}
    \,\underbrace{\widehat W_D\!\left(\frac{h}{pq}\right)}_{\text{one common physical weight}},
  \label{eq:common-weight-factor}
\end{equation}
where the argument $h/(pq)$ has no row index.  Thus the later dependence on
$m,r,p,q$ is \emph{derived dependence}; it is not freedom to choose an
independent physical weight on each row.  This distinction is load-bearing:
an arbitrary edge-dependent family would be a strictly stronger source class
and is deliberately excluded from the canonical direct theorem.

\section{Fourier normalization and determinant shell}
The normalized pre-square source is
\begin{equation}
 \widetilde C_{r,k,m}
   =\sum_{p\asymp L}b_p
    \sum_{\substack{q\asymp L\\q\ne p}}d_q
    \sum_h\omega_{m,p,q}(h)
    \e_{mq}\!\left(2hk\overline{p(m+kr)}\right),
 \label{eq:Ctilde}
\end{equation}
where
\begin{equation}
   \omega_{m,p,q}(h)
   =\frac{H}{pq}\widehat W_D\!\left(\frac{h}{pq}\right)
    \e\!\left(-\frac{h\alpha}{pqm}\right).
   \label{eq:omega}
\end{equation}
The source norms satisfy
\[
   \sum_h|\omega_{m,p,q}(h)|^2\ll H,\qquad
   \sum_h|\omega_{m,p,q}(h)|\ll H.
\]

Writing $m'=m+rk$, additive reciprocity separates the phase:
\begin{equation}
 \e_{mq}\!\left(2hk\overline{pm'}\right)
 =
 \e_{pr}\!\left(-2h\overline{qm}\right)
 \e\!\left(\frac{2h}{prqm}\right)
 \e_q\!\left(-2h\overline{prm'}\right).
 \label{eq:addrec}
\end{equation}
Thus the exact $H$-scale smooth profile is
\[
 \beta_{m,r;p,q}(h/H)
 =\frac{H}{pq}\widehat W_D(h/pq)
  \e\!\left(-\frac{h\alpha}{pqm}+\frac{2h}{prqm}\right).
\]
The $h$-Poisson/Gauss completion gives the determinant shell
\begin{equation}
    ms-prn=2.
    \label{eq:shell}
\end{equation}
Choose $\kappa_p(m)\in[0,pr)$ with $m\kappa_p(m)\equiv2\pmod{pr}$ and set
\[
    n_0(m,p,r)=\frac{m\kappa_p(m)-2}{pr}.
\]
Then every solution is parameterized by
\[
    s_j=\kappa_p(m)+jpr,\qquad n_j=n_0+jm.
\]
The quotient amplitude is
\begin{equation}
   \mathcal W_{m,r;p,q}(j)
   :=\widehat\beta_{m,r;p,q}
      \!\left(\frac{H(n_0+jm)}{qm}\right),
   \label{eq:Wsrc}
\end{equation}
and, for $|j|\ll U=L/H$,
\begin{equation}
   \sum_{|j|\ll U}|\mathcal W_{m,r;p,q}(j)|^2\ll U X^{o(1)}.
   \label{eq:Wsrcnorm}
\end{equation}

\section{Source-preserving Poisson--Bruhat transport}
\label{sec:PB}
The following source-specific statement is isolated because it is the exact
bridge between the authoritative physical source and the finite BPP compiler.
The scale identity in \cref{lem:PB-budget} is elementary; the uniform smooth
envelope, the negligible-state disposal, and the numerical source ceiling are
source/analytic inputs.

For a transported state family put
\[
 a_{r,\xi}:=|c_r(\xi)|,
 \qquad
 \mathcal T_{\rm abs}:=
 \sum_{\substack{r\in[R,2R]\\ r\ \mathrm{prime}}}
      \left(\sum_{\xi} a_{r,\xi}\right)^2.
\]

\begin{hypothesis}[One-sided source transport and absolute ceiling]
\label{hyp:PB}
For every fixed $A_0>0$, the transported source splits into a retained finite
state set $\mathfrak X_{\rm ret}$ and a discarded remainder with Gate-energy
contribution $O_{A_0}(X^{-A_0})$ times the natural scale.  Every retained state
has scalar moving-row amplitude
\begin{equation}
   c_r(\xi)=B_{\xi}F_{\xi}(r/R),
   \label{eq:smooth-envelope}
\end{equation}
where $B_{\xi}$ is independent of $r$, $F_{\xi}$ is supported in a fixed compact
subinterval of $(1,2)$, and for every fixed $J$,
\[
   \sup_{\xi}\sup_t |\partial_t^J F_{\xi}(t)|\ll_J X^{o(1)}.
\]
For every fixed truncation exponent $A$, the retained set is chosen so that
$\|F_{\xi}\|_\infty\ge 10R^{-A}$; states failing this inequality are part of
the discarded source remainder.  The transport is one-sided, has $X^{o(1)}$
nuclear cost, and gives the explicit absolute source ceiling
\begin{equation}
   \boxed{\mathcal T_{\rm abs}\ll M^2L^4X^{o(1)}.}
   \label{eq:Tabs-ceiling}
\end{equation}
No two-sided Hilbert--Schmidt equivalence is assumed.  Finally, the physical
and normalized sources obey
\[
   C_e=H^{-1}\widetilde C_e+\mathcal R_e,
   \qquad
   \sum_{e\in\mathcal E^\sharp}|\mathcal R_e|^2
     \ll_{A_0} X^{-A_0}\frac{ML^4}{H},
\]
for every fixed $A_0>0$.
\end{hypothesis}

\begin{status}[Source status]
The finite Poisson--Bruhat algebra and the budget identity below are banked.
The full assertion in \cref{hyp:PB}, including the literal source equality,
uniform retained-state truncation, the physical-normalized bridge and \eqref{eq:Tabs-ceiling}, is an explicit
source/analytic pin.  It is not inferred from the exponent arithmetic.
\end{status}

\begin{lemma}[Poisson--Bruhat budget identity]
\label{lem:PB-budget}
The scale identities imply
\[
   \boxed{\frac{MK(MH)^2}{RL^4}=1.}
\]
\end{lemma}
\begin{proof}
By \eqref{eq:scale-identities}, $MK=L^2/H$.  Therefore
\[
 \frac{MK(MH)^2}{RL^4}
 =\frac{(L^2/H)M^2H^2}{RL^4}
 =\frac{M^2H}{RL^2}=1.
\]
\end{proof}

\begin{remark}
The amplitude is smooth in $r$ but is not constant.  Thus a count of the prime
divisors of a fixed determinant does not by itself give an $R^{-1}$ weighted
family saving.
\end{remark}

\section{Band-limited prime participation}
\label{sec:bpp}
From now on $\xi$ ranges over the retained state set of \cref{hyp:PB}.  Put
$a_{r,\xi}=0$ when the prime row $r$ is excluded for the state $\xi$, and define
\[
   M_{\xi}:=|B_{\xi}|\sup_{1\le t\le2}|F_{\xi}(t)|.
\]
Fix a large $A$ and put $N=R^{1/4}$.  A smooth periodic extension of $F_{\xi}$
and repeated integration by parts give, uniformly in retained $\xi$,
\begin{equation}
   F_{\xi}(t)=P_{\xi}(t)+O_A(R^{-A}),
   \label{eq:trunc}
\end{equation}
where $P_{\xi}$ is a trigonometric polynomial of degree at most $N$.

\begin{lemma}[Bernstein plateau]
\label{lem:bernstein-plateau}
For every retained state $\xi$, there is an interval $I_{\xi}\subset[R,2R]$ of
length $\gg R^{3/4}$ on which $|c_r(\xi)|\gg M_{\xi}$.
\end{lemma}
\begin{proof}
Let $t_{\xi}$ maximize $|P_{\xi}|$.  Bernstein's inequality
\cite{Zygmund2002} gives
$\|P_{\xi}'\|_\infty\le N\|P_{\xi}\|_\infty$.  Hence on a one-sided or two-sided
interval of $t$-length $\gg N^{-1}$ contained in $[1,2]$, $P_{\xi}$ retains a
fixed proportion of its maximum.  The retained-state clause in
\cref{hyp:PB} makes the error in \eqref{eq:trunc} smaller than that maximum.
Multiplying the interval length by $R$ gives $R/N=R^{3/4}$.
\end{proof}

\begin{hypothesis}[Fixed-state exclusion]
\label{hyp:fixed-state-exclusion}
For each retained state $\xi$, every prime row excluded from the generic
non-projective branch divides one of $X^{o(1)}$ fixed nonzero integers
$N_{\xi,j}$ of polynomial height, $|N_{\xi,j}|\le X^{O(1)}$.
\end{hypothesis}

\begin{lemma}[Admissible prime participation]
\label{lem:prime-participation}
Assume \cref{hyp:fixed-state-exclusion}.  For every retained state $\xi$, the
plateau $I_{\xi}$ contains $R^{3/4-o(1)}$ admissible prime rows.  In particular,
if
\[
 A_{\xi}:=\max_{\substack{r\sim R\\r\ \mathrm{prime}}}a_{r,\xi},
 \qquad
 N_{\xi}:=\frac{\sum_{r\sim R,\ r\ \mathrm{prime}}a_{r,\xi}}{A_{\xi}},
\]
then $A_{\xi}>0$ and
\begin{equation}
   A_{\xi}\asymp M_{\xi},\qquad N_{\xi}\gg R^{3/4-o(1)}.
   \label{eq:participation}
\end{equation}
\end{lemma}
\begin{proof}
Huxley's short-interval prime number theorem applies to the interval from
\cref{lem:bernstein-plateau}, since $3/4>7/12+\varepsilon$ for a fixed small
$\varepsilon>0$ \cite{Huxley1972}.  Thus $I_{\xi}$ contains
$R^{3/4-o(1)}$ primes.  A nonzero polynomial-height integer has $O(\log X)$
distinct prime divisors, and only $X^{o(1)}$ such integers occur, so the
exclusion removes only $X^{o(1)}$ primes.  On every remaining plateau prime,
$a_{r,\xi}\gg M_{\xi}$, proving both claims.
\end{proof}

Let $\mathcal R$ be the primes in $[R,2R]$, let
$S=|\mathcal R|=R^{1+o(1)}$, and retain the notation $\mathcal T_{\rm abs}$
from \cref{sec:PB}.

\begin{lemma}[Participation-to-envelope]
\label{lem:participation-envelope}
If $N_{\xi}\ge P$ for every retained state, then
\[
   \left(\sum_{\xi}A_{\xi}\right)^2
   \le \frac{S}{P^2}\mathcal T_{\rm abs}.
\]
\end{lemma}
\begin{proof}
By Cauchy--Schwarz,
\[
 \left(\sum_{\xi}A_{\xi}N_{\xi}\right)^2
 =\left(\sum_{r\in\mathcal R}\sum_{\xi}a_{r,\xi}\right)^2
 \le S\mathcal T_{\rm abs}.
\]
Since $N_{\xi}\ge P$, the result follows.
\end{proof}

\subsection{Projective versus non-projective Gram pairs}
The transport attaches to every retained state a pair $(Z_{\xi},L_{\xi})$ and the
outer determinant
\begin{equation}
   \Delta(\xi,\eta):=Z_{\xi} L_{\eta}-Z_{\eta} L_{\xi}.
   \label{eq:outer-delta}
\end{equation}
We split the Gram pair set before applying BPP:
\[
   \Delta(\xi,\eta)\ne0\qquad\text{or}\qquad\Delta(\xi,\eta)=0.
\]
The second class is the projective branch in \cref{sec:projective}.  Since
$\Delta(\xi,\xi)=0$, the entire state diagonal lies in that branch.

For $\Delta(\xi,\eta)\ne0$, a row prime can couple the fixed pair only if it
divides the fixed nonzero polynomial-height integer $\Delta(\xi,\eta)$.  Hence
\begin{equation}
   \mathcal E_{\rm np}
   \ll X^{o(1)}\left(\sum_{\xi}A_{\xi}\right)^2.
   \label{eq:nonproj-codegree}
\end{equation}

\begin{proposition}[BPP family-energy contraction]
\label{prop:bpp}
Under \cref{hyp:PB,hyp:fixed-state-exclusion},
\[
   \boxed{\mathcal E_{\rm np}
          \ll R^{-1/2+o(1)}\mathcal T_{\rm abs}.}
\]
In particular, by \eqref{eq:Tabs-ceiling},
\[
   \mathcal E_{\rm np}
   \ll R^{-1/2+o(1)}M^2L^4.
\]
\end{proposition}
\begin{proof}
Use \cref{lem:participation-envelope} with
$P=R^{3/4-o(1)}$ and $S=R^{1+o(1)}$ in
\eqref{eq:nonproj-codegree}.
\end{proof}

\section{The unique outer square root}
\label{sec:root}
The finite outer four-cycle inequality gives a single square root, but its
application to the canonical Gate energy requires an exact energy pin.  We
state that pin explicitly.

\begin{hypothesis}[Outer-root energy pin]
\label{hyp:outer-pin}
Let $\mathcal G_{\rm np}$ denote the normalized Gate contribution of the
non-projective transported branch.  The source identification and the finite
outer-four-cycle compiler give
\begin{equation}
   \mathcal G_{\rm np}
   \ll M^2L^4X^{o(1)}
      \left(\frac{\mathcal E_{\rm np}}{\mathcal T_{\rm abs}}\right)^{1/2}.
   \label{eq:outer-pin}
\end{equation}
\end{hypothesis}

The finite square-root algebra behind \eqref{eq:outer-pin} is
\[
 \sum_r\|A_r\|_{\rm op}
 \le |\mathcal R|^{1/2}\left(\sum_r\mathcal T_r\right)^{1/2},
\]
and is machine-checked.  The equality identifying its abstract energy with
the canonical Gate energy is the source pin in \cref{hyp:outer-pin}.
Combining \cref{prop:bpp,hyp:outer-pin} gives
\[
   \mathcal G_{\rm np}
   \ll R^{-1/4+o(1)}M^2L^4.
\]
The required comparison is $MR^{-1/4}\le H$, equivalently
\[
   \frac54a+2b-1\ge0.
\]
At $V_1,V_2,V_3$ the margins are respectively
\[
   \frac1{72},\qquad\frac1{24},\qquad\frac1{32}.
\]

\section{Projective reassembly: ratio classes}
\label{sec:projective}
The complementary branch is $\Delta(\xi,\eta)=0$, including $\xi=\eta$.  On $L\ne0$
the correct class is the projective ratio $\kappa(Z,L)=Z/L$, not the product
$ZL$.  For a finite transported state support $\mathcal S$, put
$v_{Z,L}=A_Z\otimes B_L$ and
\[
   C_\xi:=\sum_{(Z,L)\in\mathcal S:\ \kappa(Z,L)=\xi}v_{Z,L}.
\]
Then the exact class-pushforward identity is
\begin{equation}
 \sum_{\xi}\|C_\xi\|^2
 =\sum_{(Z_1,L_1),(Z_2,L_2)\in\mathcal S}
   \1_{Z_1L_2=Z_2L_1}
   \operatorname{Re}\langle v_{Z_1,L_1},v_{Z_2,L_2}\rangle.
 \label{eq:ratio-energy}
\end{equation}

\begin{hypothesis}[Projective source routing]
\label{hyp:projective}
The transported $\Delta=0$ source is routed into positive primitive
ratio-class packets of class multiplicity $X^{o(1)}$.  One-zero axes and
non-primitive/rank-loss states are routed to explicit exceptional blocks.
Moreover the diagonal/projective source mass satisfies
\begin{equation}
   \sum_{(Z,L)\in\mathcal S}\|A_Z\|^2\|B_L\|^2
   \ll MHL^4X^{o(1)},
   \label{eq:projective-source-ceiling}
\end{equation}
and every exceptional block is bounded at the same target.
\end{hypothesis}

Under \cref{hyp:projective}, Cauchy within each ratio class and
\eqref{eq:ratio-energy} give
\[
   \mathcal G_{\rm proj}\ll MHL^4X^{o(1)}.
\]
The finite ratio-class identity, primitive rigidity and multiplicity lemma are
machine-checked; \cref{hyp:projective} is the physical source-routing pin.

\section{Recombination error}
\begin{hypothesis}[Literal source recombination estimate]
\label{hyp:recombination}
The quotient-cell/source recombination error obeys
\begin{equation}
   \mathcal E_{\rm rec}
   \ll U^{-2+o(1)}M^2L^4.
   \label{eq:recombination-error}
\end{equation}
\end{hypothesis}
Since $U=L/H$ and $DH=L^2$,
\[
 \frac{U^{-2}M^2L^4}{MHL^4}=\frac{M}{D}.
\]
The margins at $V_1,V_2,V_3$ are $1/18,1/18,1/24$.

\section{Conditional main theorem}
\begin{theorem}[Canonical all-$m$ common-weight estimate]
\label{thm:main}
Let the parameters lie in the polytope generated by $V_1,V_2,V_3$.  Let
$\mathcal E^\sharp$ be the generic unit/full-conductor all-$m$ row family and
let $C_e(b,d)$ be the covariance \eqref{eq:Ce}.  Assume
\cref{hyp:PB,hyp:fixed-state-exclusion,hyp:outer-pin,hyp:projective,hyp:recombination}.
Then
\[
   \boxed{\sum_{e\in\mathcal E^\sharp}|C_e(b,d)|^2
          \ll \frac{ML^4}{H}X^{o(1)}.}
\]
Equivalently,
\[
   \boxed{\sum_{e\in\mathcal E^\sharp}|\widetilde C_e(b,d)|^2
          \ll MHL^4X^{o(1)}.}
\]
\end{theorem}
\begin{proof}
The non-projective branch is bounded by
\cref{prop:bpp,hyp:outer-pin}.  The complete $\Delta=0$ branch, including the
state diagonal, is bounded by \cref{hyp:projective} and the exact ratio-class
identity.  The discarded source and the recombination error are controlled by
\cref{hyp:PB,hyp:recombination}.  These branches are exhaustive and disjoint.
The normalized estimate follows.  The physical-normalized identity in
\cref{hyp:PB}, together with its square-summable remainder, gives the physical
estimate.
\end{proof}

\begin{corollary}[Clean-$P_3$ restriction]
Every clean-$P_3$ subfamily of $\mathcal E^\sharp$ satisfies the same bound.
\end{corollary}
\begin{proof}
The energy is a sum of non-negative row energies.
\end{proof}

\section{Scope separation: Gate 1A versus source routing}
The theorem above is a conditional analytic theorem for the canonical direct operator.
It does not assert that every packet in a global sieve decomposition has
already been identified with this operator.

Suppose an upstream source compiler gives an exact decomposition
\[
  \mathcal S_{\mathrm{actual}}
  =\sum_{\lambda\in\Lambda_{\mathrm{gen}}}
      a_\lambda D_\lambda U_\lambda\mathcal S_\lambda^\sharp
    +\mathcal S_{\mathrm{exc}},
\]
where $U_\lambda$ is unitary/permutational and
\[
   \sum_\lambda |a_\lambda|\|D_\lambda\|_{\mathrm{op}}\le X^{o(1)}.
\]
Then the main theorem controls the generic packets by Minkowski, provided each
packet is truly an instance of the canonical operator. Proving the source
identity, packet multiplicity, and exceptional routing is a separate Gate-0
compiler problem.

In particular, a common physical weight routes directly into the theorem.
A fixed-dimensional smooth family generated deterministically from the common
source can be decomposed into $X^{o(1)}$ common templates.  By contrast, an
arbitrary independently varying edge weight is a genuinely stronger object.
In the formal development this more general datum is retained as an upstream
Gate-0 source-adapter obligation; it is not a premise of the canonical
Gate-1A direct theorem.

This separation matters downstream. The Ford--Maynard framework applies to
the actual non-negative sequence and its complete Type-I/Type-II information
\cite{FordMaynard2024}; a convenient analytic model cannot replace an
unverified source census.

\section{Formal verification and provenance}
\label{sec:formal-status}

The accompanying Lean bank deliberately separates finite mathematics, source
provenance, and deep analytic input.  This is useful here because a theorem
compiler can be perfectly correct while its external analytic hypothesis is
not itself formalized in Mathlib.

\subsection{Machine-checked finite/source skeleton}
The v9.8 formal bank verifies, without user axioms or admitted proofs:
\begin{itemize}
  \item the canonical pre-square direct source as a literal source
        transcription, together with the fact that the physical transform
        $\widehat W_D$ is one common datum;
  \item the all-$m$ generic row type, with no clean-$P_3$ field;
  \item the exact normalized and physical Gate-1A energies and their
        $H^2$ normalization bridge;
  \item the finite participation-to-family-energy inequality;
  \item the exponent conversion
        $R^{-1/2}\mapsto R^{-1/4}$ and the margins
        $1/72,1/24,1/32$;
  \item the projective/axis finite algebra and the $U^{-2}$ exponent budget;
  \item the logical separation between the Gate-1A direct theorem and the
        upstream global packet/source census.
\end{itemize}

The relevant files are
\begin{center}
\begin{tabular}{>{\ttfamily}l l}
\toprule
V98CanonicalDirectSource.lean & canonical source and common-weight firewall\\
V98CanonicalAllMRows.lean & all-$m$ row family\\
V98DirectEnergyPin.lean & exact energy objects and pinning interface\\
V98BPPProvenance.lean & external BPP participation provenance\\
V98DirectClosure.lean & deterministic normalized-to-physical compiler\\
V98Gate0ScopeSplit.lean & Gate-0/Gate-1A scope separation\\
V98Status.lean & axiom/status audit\\
\bottomrule
\end{tabular}
\end{center}

\subsection{What Lean does not prove}
The formal development does \emph{not} manufacture an inhabitant of the deep
analytic BPP source input.  In particular, the trigonometric Bernstein
estimate, the source-specific smooth truncation, and Huxley's short-interval
prime theorem are not currently formalized in the project.  They are recorded
with external/source-specific provenance, and no global axiom is introduced.
The state-diagonal issue is not external: the finite moving-family bank splits
diagonal and off-diagonal exactly, and the corrected projective identity places
the diagonal in the $\Delta=0$ branch.

Likewise, the formal ``energy pin'' is an explicit equality/interface tying a
generic BPP energy to the canonical Gate-1A energy; the development includes
guards showing that an arbitrary abstract bound cannot manufacture this pin.
Thus a successful compiler is not confused with an independently proved
analytic estimate.

\subsection{Research status versus formal status}
Accordingly the appropriate status language is:
\[
 \boxed{\text{source-explicit conditional theorem under the displayed analytic/source pins}}
\]
for the mathematics of this note, and
\[
 \boxed{\text{external-uninhabited in Lean}}
\]
for the fully formal end-to-end analytic proof.  The latter is a provenance
boundary, not an exponent deficit or an additional arithmetic obstruction.



\subsection{Gate-1B v8.5 H7 short-short bank}
The append-only v8.5 bank separates the H7 short-short region from the
high-prime complement.  It machine-checks the finite common-sequence and
joint-prime compiler and records the exact capacity conversion
\[
     Y^{-1/2}=X^{-1/18}.
\]
Its closure theorem is explicitly conditional on the source/common-sequence,
source-energy and multiplicative-large-sieve inputs; no unconditional
\texttt{H7\_CLOSED} theorem is exported.  The same bank leaves the high-prime
complement open and records H8 only as an uninhabited robustness interface.
This formal provenance is why the stronger H8 statements elsewhere in the
draft are labelled research-capacity audits rather than banked closure
theorems.

\subsection{Gate-1B V10 zero/source/reassembly bank}
The latest append-only Gate-1B V10 bank adds an exact finite/algebraic zero-mode
and reassembly layer without altering the earlier Gate-1A bank.  In particular,
it reproves the canonical additive-character baseline from the standard
character system, proves that the canonical discrepancy has no additive
zero-frequency, and proves the historical/canonical residual identity
\[
   S_{\mathrm{hist}}=S_{\mathrm{can}}-R_E.
\]
A perturbation counterguard shows that changing an abstract historical
comparison $E(q)$ changes $R_E$ while leaving the nonzero-frequency packet
unchanged.  Thus the canonical zero mode is closed at the finite/algebraic
level, whereas identifying a historical $E(q)$ with a canonical comparison is
an independent source question.

The same V10 bank proves literal finite packet partitions and deterministic
packet-budget/reassembly compilers, including conditional implications from
exact leaf inputs to the project's local Gate-1B and Type-II predicates.  It
does \emph{not} inhabit the analytic leaves.  A repository-wide audit in that
run also found no literal declaration named \texttt{FullFMTypeII}\allowbreak\texttt{\_OneSixth};
the formal compiler consequently targets the project's existing local Type-II
predicate and must not be described as a formal Ford--Maynard one-sixth theorem.
The V10 build is sorry-free and introduces no user axioms.

\section{Relation to Kloosterman-fraction technology}
Bettin--Chandee prove non-trivial bounds for separated trilinear Kloosterman
fractions with arbitrary coefficient sequences \cite{BettinChandee2015}.
Wright improves this architecture when a factor of the denominator is fixed
\cite{Wright2026}.  Blomer--Pascadi obtain a $c^{-1/32}$ saving in the
square-root critical range for fixed-modulus bilinear Kloosterman forms
\cite{BlomerPascadi2026}, while Pascadi's non-abelian amplification gives, for
example, a $c^{-1/12}$ saving for balanced products of two comparable primes
\cite{Pascadi2025}.  These are important fixed-multiplier/fixed-modulus
capacity inputs.  They do not, by their published statements, estimate the
coherent moving source-multiplier family isolated in the Gate-1B appendix.
The final Gate-1A direct proof likewise does not black-box any of these
Kloosterman theorems; its decisive resource is smooth participation in the
moving prime family.

\section{What remains outside this theorem}
The following are not conclusions of the main theorem:
\begin{enumerate}
  \item exhaustive routing of every global high-complexity source packet into
        the canonical direct operator;
  \item closure of the separate Gate-1B analytic frontier, whose present
        high-conductor residual is recorded in \cref{app:gate1b-frontier,app:hicond-nearprim,app:mam45-native,app:pure5-dp,app:rankone-bezout};
  \item Full Type-II reassembly;
  \item a Hardy--Littlewood prime-pair asymptotic or twin-prime infinitude.
\end{enumerate}

The point of the second item is structural.  The conditional Gate-1A theorem
established here supplies one all-$m$ analytic engine once its displayed pins
are verified.  The complementary programme has progressed beyond the earlier
single H7 natural-scale residual.  The high-conductor audit first isolated an
unavoidable near-primitive full-conductor signed packet and then identified its
source-minimal native parent: the global $4$M$\mid5$D MAM45 scalar with linear
$\mu(d)$ and the long prime $p$ lane.  The top near-primitive packet is an exact
child of this parent, while proper-divisor terms are triangular $r>1$ children.
Neither Gate~1A, the common-conductor square norm, nor a coefficient-blind
fixed-multiplier Kloosterman estimate supplies this global signed theorem.

\section{Conclusion}
The main conceptual point is that local sparsity alone does not give a
moving-family saving. After source-preserving transport, the row coefficient
remains $r$-dependent through a uniformly smooth envelope. Band-limited prime
participation supplies the missing family coherence: smoothness creates a long
plateau, prime distribution fills it with many moving rows, and finite Cauchy
converts participation into the required $R^{-1/2+o(1)}$ energy contraction.
After one and only one outer square root, the resulting $R^{-1/4+o(1)}$ gain
closes the canonical all-$m$ common-weight Gate-1A direct estimate throughout
the parameter polytope.

A substantive feature of this mechanism is that its resource is the moving
prime row $r$, not a prime factor of $m$.  No largest-prime-factor quantity of
$m$ enters \cref{lem:bernstein-plateau,lem:prime-participation,prop:bpp}.
Thus the argument applies uniformly to the canonical all-$m$ row family,
including smooth values of $m$; clean-$P_3$ is a positive-restriction
corollary rather than the source of the saving.


On the complementary side, Late update~VIII now supersedes the Draft~15
fully-crossed headline.  The entropy trichotomy and prime-divisor rigidity
survive, but the central band has been reduced further by two GCD extractions
to a B\'ezout line.  Large $gs$ tails are disposable with arbitrary logarithmic
saving under the banked fixed-polylog multiplicity, while the primitive
$g=s=1$ line still has the full natural $QV$ scale.  Hence
\texttt{FULLCROSS-RIGIDITY-ONLY-CLOSURE45} is a banked failed route: support,
divisor multiplicity and injectivity can sparsify the incidence graph but
cannot create cancellation.

The new canonical primitive child is obtained from
\[
 \beta_{D,P}(n)=-\mu(n)L_{D,P}(n)
\]
and the determinant-$2$ B\'ezout orbit
\[
 \ell A_0-u w_0=2,\qquad A_t=A_0+ut,\qquad w_t=w_0+\ell t.
\]
It is the source-specific rank-one correlation
\texttt{RANKONE-BEZOUT-AFFINE-ELLIOTT5D45}.  The shift pair
$(uk,\ell k)$ ranges over only an $H^{-1}$-density subset of the independent
shift rectangle, with $H\ge X^{5/18}$, so complete-shift averaged-Chowla
theorems do not directly pay for this restriction.  Conversely the present
extra averages and determinant constraint mean that the new theorem has not
been shown equivalent to ordinary fixed-shift Chowla.  The low pretentious
profile is internally eliminated; the genuine first unproved implication is
an inverse statement forcing either a common low profile or a quantitatively
controlled high-conductor/high-frequency packet, with enough entropy penalty
to obtain $QV(\log X)^{-A}$.  The best presently banked scale is only
$QV(\log X)^{C_0}$, so the remaining deficit is logarithmic rather than a
power of $X$.

A source-preserving reverse audit also recovers MRK5 and manufactures the
MOYK5 Kloosterman kernel, but Yang-type factorable-weight landing is false,
the Blomer--Pascadi global application is power-nonclosing in the physical
lengths, a second transform returns the QK5 representation, and Wright's
partially fixed-denominator theorem remains at least $X^{1/2-o(1)}$ above the
physical scalar target in the most favourable split.  These are death
certificates against representation shopping, not evidence that the
rank-one theorem is impossible.  The comparison/main-term pin, lower-defect
packets, the $r>1$ square-character family, transition strip, QK56
parent/exhaustiveness and literal source reassembly remain open obligations.

Downstream, Ford--Maynard's universal Type-II hypothesis remains the published
sufficient wrapper.  Their proof of Proposition~7.22 identifies equation~(7.23)
as a narrower proof-specific landing point, while the latest source archaeology
suggests that the remaining Gate-to-Ford problem is first a source-exact packet
dictionary rather than another abstract rank-reduction theorem.  We label that
interface \texttt{SOURCE-ADAPTER45}.  Thus the uncertainty has narrowed
substantially, but Gate~1B, the downstream source dictionary, the Ford--Maynard
endgame, and twin-prime infinitude all remain open/conditional as stated.


\appendix
\section{Formal theorem map and scope ledger}
\label{app:formal-map}
\begin{center}
\begin{tabular}{p{0.31\textwidth} p{0.23\textwidth} p{0.36\textwidth}}
\toprule
Component & Status & Role in this note\\
\midrule
Canonical source and common $W_D$ & source-pinned; finite identities checked & fixes the source class\\
All-$m$ row family & Lean-constructed & removes clean-$P_3$ from the statement\\
PB smooth envelope, retained-state disposal and $\mathcal T_{\rm abs}$ ceiling & \textbf{conditional source/analytic pin} & supplies the absolute comparator\\
Finite BPP compiler & Lean-proved & participation $\Rightarrow$ family energy\\
Huxley + fixed-state exclusion & published theorem + source pin & supplies prime participation\\
Outer-root energy pin & \textbf{conditional source pin}; square-root algebra proved & identifies the abstract energy with the Gate energy\\
Ratio-class projective algebra & Lean-proved finite algebra & includes the diagonal in $\Delta=0$\\
Physical projective routing and source ceiling & \textbf{conditional source pin} & bounds the zero-curvature branch\\
$U^{-2}$ exponent budget & Lean-proved arithmetic & checks the recombination margin\\
Literal recombination estimate & \textbf{conditional source pin} & controls the source error\\
Gate-1B canonical zero mode & V10 Lean-proved finite algebra & zero additive frequency removed canonically\\
H7 short-short joint-prime compiler & conditional analytic compiler; $X^{-1/18}$ capacity & source/common-sequence pins remain\\
H8 analogue & research capacity audit; formal robustness interface open & not a banked closure theorem\\
Same-$q$ QK56 cross-character Gram & \textbf{internal analytic pass on proper smooth sector} & programme claim, not published/Lean-closed\\
Cross-$q$ source spread & \textbf{internal finite/source pass} & programme source-spread audit; final transcription pending\\
Shift $TT^*$ literal source pin & \textbf{source open} & recover the literal cyclic source before the analytic child\\
Shifted source-linked character child & historical research child; necessity now undetermined & superseded as first open by Late update~VIII\\
Arbitrary edge-dependent packet & outside theorem scope & upstream adapter obligation\\
Global packet census / Full Type II & open & not implied here\\
\bottomrule
\end{tabular}
\end{center}
\section{Canonical source versus arbitrary edge dependence}
\label{app:edge-dependence}

It is useful to make the source distinction algebraically explicit.  The
canonical coefficient has the form
\[
   \omega_e(p,q,h)
   =R_e(p,q,h)\,\widehat W_D(h/(pq)),
\]
where $R_e$ is a deterministic scalar generated by the row and by exact
reciprocity.  Hence the entire family is determined once the one physical
weight $W_D$ is fixed.  This is fundamentally different from prescribing an
arbitrary collection
\[
   \{W_{D,e}:e\in\E^\sharp\}
\]
with no relation between different rows.  The latter can have effective rank
comparable with the number of rows and therefore cannot be reduced to the
canonical theorem by functional analysis alone.  This is why the global
source compiler must either identify a packet with the common-weight source,
reduce it to finitely many smooth templates, or route it elsewhere.



\section{Complementary programme boundary: Gate 1B}
\label{app:gate1b-frontier}
This appendix is a research-status calculation, not an ingredient in
\cref{thm:main}.  ``Exact'' below refers to the displayed squarefree/unit
source dictionary and Fourier convention.  The natural-scale estimate in
\cref{prop:h7-natural} additionally uses the explicit source norm in
\cref{hyp:h7-source-norm}.

\subsection{Low/medium conductor and the high-conductor threshold}
Write $q=ce$, $c\asymp C$, $q\asymp Q$.  After the corrected state count,
\[
   \sum_{c,\chi^\ast,e}|\beta(ce)|^2\ll QC\,X^{o(1)}.
\]
The coefficient-blind estimate has the relative shape
\[
 \frac{|\mathcal H_s(C)|}{Y^9}
 \ll
 \begin{cases}
  C/Q,& C\le Y^4,\\
  C^2/(QY^4),& C\ge Y^4.
 \end{cases}
\]
The formulae cross at $C=Y^4$.  The large-$C$ estimate loses all saving at
\[
   C_0=Q^{1/2}Y^2,
\]
so $C_0$ is the zero-saving/high-conductor threshold, not the crossing point.

\subsection{Exact source split and induced-Gauss cancellation}
The squarefree coefficient is
\[
 \beta(q)=\sum_{dp=q}\mu(d)\log p,
\]
with $p>V$ and $q/p>U$.  In high conductor, $e=q/c\le Y^2$, whereas
$V=Y^{5/2-o(1)}$, hence $p>V>e$ and the lane $p\mid e$ is empty.  Thus
\[
   c=pc_0,\qquad d=c_0e,
\]
and
\begin{equation}
   \beta(ce)=\mu(e)\rho_{C,e}(c),
   \qquad
   \rho_{C,e}(c)=
   \sum_{pc_0=c}^{\rm physical}\mu(c_0)\log p.
   \label{eq:beta-C}
\end{equation}
For a primitive $\chi^\ast\pmod c$ induced to $ce$,
\begin{equation}
   \tau_{ce}(\chi)=\mu(e)\chi^\ast(e)\tau_c(\chi^\ast).
   \label{eq:induced-gauss2}
\end{equation}
Consequently the short-cofactor sign is spent exactly:
\[
   \beta(ce)\tau_{ce}(\chi)
   =\rho_{C,e}(c)\chi^\ast(e)\tau_c(\chi^\ast).
\]
It cannot be used a second time as cancellation.

\subsection{Prime-character collapse}
Write $c=pc_0$, $(p,c_0)=1$, and factor a primitive character as
$\chi=\chi_p\chi_0$.  Let
$N=r_1\cdots r_7x\asymp Y^8$.  The $p$-local argument is
\[
   A_p\equiv c_0eN\,\overline{2h}\pmod p.
\]
For $A\in\mathbb F_p^\times$,
\begin{lemma}[Prime-character collapse]
\label{lem:prime-collapse}
\[
 \sum_{\chi_p\ne\chi_{0,p}}
   \frac{\tau_p(\chi_p)}{\sqrt p}\chi_p(A)
 =\frac{p-1}{\sqrt p}e_p(\bar A)+\frac1{\sqrt p}.
\]
\end{lemma}
\begin{proof}
Summing first over all multiplicative characters gives
$(p-1)e_p(\bar A)$.  The principal character has Gauss sum $-1$, so removing
it adds $1$.
\end{proof}
The main term therefore becomes the linear additive phase
\[
   e_p\!\left(2h\,\overline{c_0eN}\right).
\]
The correction is smaller than the main local multiplier by $p^{-1}$; since
$p>V=Y^{5/2-o(1)}$, it is power-negligible under the same absolute source
ceiling used below.

\subsection{Hybrid Poisson formula and primitive projector}
Use $\widehat v(\xi)=\int v(x)e(-x\xi)\,dx$.
\begin{lemma}[Hybrid Poisson]
\label{lem:hybrid-poisson}
Let $p\nmid c_0$, and let $\chi_0$ be primitive modulo $c_0$.  With the
consistent Fourier sign convention,
\[
 \sum_{h\in\mathbb Z}v(h/H)\overline{\chi_0(h)}e_p(ah)
 =\frac{H}{c_0}\tau_{c_0}(\overline{\chi_0})
  \sum_{\substack{m\in\mathbb Z\\m\equiv ac_0\ (p)}}
   \chi_0(-m\bar p)\widehat v\!\left(\frac{mH}{pc_0}\right).
\]
\end{lemma}
\begin{proof}
Apply Poisson summation in residue classes modulo $pc_0$.  The finite Fourier
sum factors by CRT.  The $p$-factor enforces the displayed congruence, while
the $c_0$-factor is the twisted Gauss sum
$\tau_{c_0}(\overline{\chi_0})\chi_0(-m\bar p)$.  The factor $p$ from the
$p$-sum cancels the $p$ in the Poisson denominator.
\end{proof}
Here $a=2\overline{c_0eN}\pmod p$, so the surviving frequency obeys
\[
   meN\equiv2\pmod p,
   \qquad |m|\ll(Y/e)X^{o(1)}<p.
\]
For odd squarefree $c_0$ and $(A,c_0)=1$,
\begin{lemma}[Primitive-character projector]
\label{lem:primitive-projector}
\[
  \sum_{\chi_0\ (\mathrm{mod}\ c_0)}^*\chi_0(A)
  =\sum_{\substack{d\mid c_0\\d\mid A-1}}
     \mu(c_0/d)\phi(d).
\]
Consequently
\[
 \frac{\mu(c_0)}{c_0}
 \sum_{\chi_0}^*\chi_0(A)
 =\frac1{c_0}
  \sum_{\substack{d\mid c_0\\d\mid A-1}}\mu(d)\phi(d).
\]
\end{lemma}
\begin{proof}
This is Möbius inversion on the conductor of a character.  The second identity
uses squarefreeness: $\mu(c_0)\mu(c_0/d)=\mu(d)$.
\end{proof}
With $A=meN/2$, the two projectors combine into $pd\mid meN-2$.

\subsection{Dual determinant and the fixed-power repair}
The main term is reduced to
\begin{align}
 \mathcal D_7(C)=
 \sum_{\substack{pc_0\asymp C\\p>V\\e\asymp Q/C}}
 &\log p\left(1-\frac1p\right)
 \sum_{N\asymp Y^8}a_7(N)
 \sum_{\substack{|m|\ll(Y/e)X^{o(1)}\\p\mid meN-2}}
 \widehat v(me/Y) \notag\\
 &\qquad\times\frac1{c_0}
 \sum_{\substack{d\mid c_0\\d\mid meN-2}}\mu(d)\phi(d).
 \label{eq:D7}
\end{align}

\begin{hypothesis}[H7 source norm and cell multiplicity]
\label{hyp:h7-source-norm}
The literal seven-defect plus one-model coefficient satisfies
\[
   \sum_{N\asymp Y^8}|a_7(N)|\ll Y^8X^{o(1)},
\]
$\widehat v$ is rapidly decreasing with bounded supremum, and the omitted
physical dyadic/coprimality selectors have $X^{o(1)}$ total multiplicity.
The primitive character family in \cref{lem:primitive-projector} is the
literal family produced by the source dictionary (with any parity restriction
separately accounted for).
\end{hypothesis}

\begin{proposition}[Natural-scale fixed-power repair]
\label{prop:h7-natural}
Under \cref{hyp:h7-source-norm},
\[
   |\mathcal D_7(C)|\ll Y^9X^{o(1)}.
\]
\end{proposition}
\begin{proof}
For fixed $(e,m,N)$, the nonzero integer $meN-2$ has $O(1)$ prime divisors
$p>V=Y^{5/2-o(1)}$.  After $p$ is fixed, write $c_0=ds$.  The absolute divisor
projector is bounded by
\[
 \sum_{d\mid meN-2}\ \sum_{s\asymp C/(pd)}
   \frac{\phi(d)}{ds}
 \ll X^{o(1)},
\]
because the $s$-sum is harmonic on a dyadic interval and the divisor count is
$X^{o(1)}$.  Summing $e\asymp E$ and
$|m|\ll Y/e$ costs $O(YX^{o(1)})$.  Finally use the $\ell^1$ source norm of
$a_7$.
\end{proof}
Thus the former coefficient-blind fixed-power loss $Q/Y^4$ is recovered, but
the target is $Y^9(\log X)^{-A}$.  The remaining problem is therefore the
arbitrary-log estimate
\[
   \boxed{\mathcal D_7(C)\ll_A Y^9(\log X)^{-A}.}
\]
Equivalently, one may formulate a centred seven-defect determinant variance
for
\[
   nN-pd\ell'=2,
   \qquad n\ll Y,
   \qquad N\asymp Y^8.
\]
This estimate is open.  The transform also has an anti-loop: the full divisor
term $d=c_0$ reconstructs the original two-model determinant skeleton, so the
Poisson/character transform is not an iteration-to-closure by itself.

\subsection{Historical H7 natural-scale status and the later promotion}
\label{sec:h7-historical-promotion}
The natural-scale proposition above records an important historical reduction:
it removes the former fixed-power high-conductor deficit.  It is no longer the
controlling final status for the entire H7 analysis.  In the short-short region
\[
   \alpha<\frac49,\qquad \beta<\frac49,
\]
the v8.5 bank isolates a joint-prime character packet and proves a deterministic
multiplicative-large-sieve compiler.  Under the explicit common-sequence,
source-energy and multiplicative-large-sieve inputs, the resulting target has
relative capacity
\[
    Y^{-1/2}=X^{-1/18}.
\]
This is a \emph{conditional analytic compiler}, not an unconditional H7
closure theorem.  The high-prime complement is disjoint from this region and
remains a separate source/routing problem in the formal v8.5 bank.

The latest research audit further indicates that the H8 analogue has the same
fixed-power capacity after the corresponding source-energy audit.  We record
that as research-level analytic capacity only: the v8.5 formal development
still leaves the H8 robustness interface uninhabited.  Thus the publication
status is deliberately asymmetric:
\[
\begin{aligned}
 \text{H7 short-short:}&\quad \text{conditional compiler with }X^{-1/18}\text{ capacity},\\
 \text{H8:}&\quad \text{research capacity audit; formal source/robustness pin open}.
\end{aligned}
\]

\subsection{The full \texorpdfstring{$4/9$}{4/9} complement and the square-root transition}
The late source census organises the complement by prime exponents.  At the
programme level one separates a high-$d$ wing and a high-$p$ wing.  Joint-prime
machinery has favourable capacity on the high-$d$ side and on the high-$p$
side away from the square-root transition.  The strip in which the large
prime approaches $X^{1/2}$ is not absorbed by the short-short H7 compiler and
is instead placed in the shifted-quotient parent introduced below.

This subsection is a routing statement, not a completed analytic theorem.  In
particular, the v8.5 formal bank still records the high-prime complement as
open.  The purpose of the parent formulation is to avoid pretending that a
bound on one exponent cell automatically transports to a disjoint cell.

\subsection{The Full-Nine shifted-quotient parent}
The high-prime switch and the pure H9 specialization are most economically
viewed inside one shifted-quotient parent of the schematic form
\begin{equation}
   \mathcal S_{\mathrm{SQ}}
   =\sum_{mp\asymp X}\gamma(m)\lambda(p)\log p\;c_{\mathrm{res}}(mp-2).
   \label{eq:fullnine-shifted-quotient}
\end{equation}
Here $\gamma$ and $\lambda$ retain the actual source-generated convolution and
sign structure.  The high-prime $M$-switch is a subrange of
\eqref{eq:fullnine-shifted-quotient}, while H9 is a pure-nine-defect
specialization.  Type-I and joint-large-sieve mechanisms are plausible on the
strongly unbalanced sides; the balanced square-root strip is the hard
transition.

Equation~\eqref{eq:fullnine-shifted-quotient} is a programme parent operator,
not a theorem asserted closed in this manuscript.

\subsection{The QK56 full-covariance parent}
A second parent keeps the complete covariance of the one-copy QK output.  If
$\mathcal Q(q;h,k)$ denotes the one-copy output after the exact local
transforms, the parent is schematically
\begin{equation}
   \mathcal C_{\mathrm{QK}}
   =\sum_{q_1,q_2}\sum_{h_1,k_1,h_2,k_2}
      A(q_1,q_2;h_1,k_1,h_2,k_2)\,
      \mathcal Q(q_1;h_1,k_1)\overline{\mathcal Q(q_2;h_2,k_2)}.
   \label{eq:qk56-covariance}
\end{equation}
The decomposition
\[
   q_1=q_2 \quad\text{versus}\quad q_1\ne q_2
\]
is precisely the same-$q$ diagonal versus cross-modulus off-diagonal of this
single covariance parent.  The same-$q$ object does not collapse to a scalar
residue-energy theorem in the source analysis; the surviving object is a
multi-factor character Gram.  The cross-modulus term retains the moving
source data and likewise cannot be replaced by an arbitrary fixed-modulus
coefficient sequence.

\subsection{Pair-modulus capacity and the moving source-multiplier frontier}
Both the balanced shifted-quotient $TT^*$ child and the QK cross-modulus child
lead to product moduli.  For a \emph{fixed} multiplier, current published
Kloosterman technology has genuine power capacity: Blomer--Pascadi obtain a
$c^{-1/32}$ saving in the square-root critical range
\cite{BlomerPascadi2026}, and Pascadi's non-abelian amplification gives, for
balanced products of two comparable primes, a $c^{-1/12}$ saving
\cite{Pascadi2025}.  These results justify a backend-capacity statement only.

The actual parent carries a coherent source family
\begin{equation}
    \sum_{\Theta} A_c(\Theta)\,\mathcal K_{c,\Theta},
    \label{eq:moving-source-multiplier}
\end{equation}
where the multiplier $A_c(\Theta)$ varies with the physical packet.  A theorem
for each fixed $\Theta$ does not by itself bound
\eqref{eq:moving-source-multiplier}.  This pair-modulus formulation remains a
useful historical umbrella.  The later audit recorded in
\cref{app:late-update-II} splits the QK56 covariance into same-$q$ and cross-$q$
children and moves the principal analytic candidate to a source-linked shifted
child.  Accordingly the labels
\[
   \texttt{SOURCE-QUADRATIC-MULTIPLIER45},\qquad
   \texttt{FM-PERRON-PAIRMOD-SOURCE-MULT45}
\]
are retained only as umbrella/source-grammar labels, not as the controlling
first-open theorem of Draft~10.

\subsection{Updated Gate-1B ledger}
The current publication-facing ledger is:
\begin{center}
\begin{tabular}{p{0.43\textwidth}p{0.45\textwidth}}
\toprule
Component & Current status\\
\midrule
Orders $1$--$4$ / low-order finite geometry & banked structural material\\
H6 regroup geometry & banked structural material\\
H7 short-short & conditional joint-prime compiler; $X^{-1/18}$ capacity\\
H8 short-short analogue & research capacity audit; formal robustness/source pin open\\
High-$d$/high-$p$ complement & routing/capacity programme; square-root strip open\\
Full-Nine shifted-quotient parent & research parent; source-linked shifted child remains open\\
QK56 high-conductor covariance & \textbf{OPEN; common-conductor block is secondary after near-primitive audit}\\
Canonical additive zero mode & V10 exact finite/Lean pass\\
Historical $E(q)$ residual & source-identification fork if historical formulation retained\\
Same-$q$ QK56 cross-character proper-smooth child & \textbf{internal analytic pass; not a standalone high-conductor sign closure}\\
Coprime cross-$q$ no-wrap operator & \textbf{internal operator pass; high-conductor signed source family still open}\\
High-conductor signed parent & \textbf{RECOVERED internally}\\
Common-conductor signed covariance & exact square norm; \textbf{secondary $r>1$ child, not full parity parent}\\
Old SHORT-RJ pre-Cauchy dictionary & \textbf{FALSE / RETRACTED}\\
Square-complement Pascadi route & \textbf{retired: naive dictionary false; legal lanes power-nonclosing}\\
Shift $TT^*$ literal source pin & \textbf{SOURCE OPEN}\\
Near-primitive full-$c$ five-defect signed child & \textbf{FIRST HIGH-CONDUCTOR ANALYTIC OPEN}\\
Shifted source-linked character child & source/analytic open, downstream of literal source pin\\
S1/S2 literal source pins & narrow source interfaces remain\\
Packet census / Gate-1B reassembly & open\\
Gate 1B & \textbf{NOT CLOSED}\\
\bottomrule
\end{tabular}
\end{center}

The older D12 bulk/spike/Pascadi shortcut remains retired; this does not erase
the underlying cross-modulus contribution, which is now represented in the
covariance/pair-modulus parent.

\section{Gate 0, Gate 2, and the Ford--Maynard endgame}
\label{sec:gate02}

The preceding sections isolate the two Type-II analytic engines.  To make the
programme as close to a complete conditional theorem as the present evidence
permits, we now state the remaining Gate-0/Gate-2 interface explicitly.  This
section is downstream bookkeeping and source analysis; it is not used in the
proof of the conditional Gate-1A theorem.

\subsection{Gate 0: deterministic Type-I compiler and research status}
The separate Gate-0 bank works with a shifted-prime sequence $a_n$, a simpler
comparison sequence $b_n$, and $w_n=a_n-b_n$.  Its deterministic compiler has
the following logical form.  A suitable residue-class weighted maximal
Bombieri--Vinogradov input for $a_n$, together with the comparison progression
mean for $b_n$ and the exact reindexing bridge, implies the Ford--Maynard Type-I
predicate for $w_n$.  The finite compiler is machine-checked, while the deep
analytic inputs remain external to Lean.

Accordingly we record the Gate-0 status as
\[
 \boxed{\text{research-level analytic pass; external-uninhabited in Lean}.}
\]
This status does not contain any Type-II information and therefore cannot by
itself activate the Ford--Maynard lower-bound theorem.

\subsection{The operating Ford--Maynard parameters}
For the downstream application we choose the central operating width
\[
    \nu_0=\frac16.
\]
For $\varepsilon>0$ set
\begin{equation}
  P_\varepsilon
   =\left(\frac12-\varepsilon,\,\varepsilon,\,
            \frac16-2\varepsilon\right).
  \label{eq:Peps}
\end{equation}
The corresponding Type-II interval in the Ford--Maynard convention is
\begin{equation}
    [\theta,\theta+\nu]
    =\left[\varepsilon,\frac16-\varepsilon\right].
    \label{eq:fm-typeii-window}
\end{equation}
The published central threshold in Ford--Maynard Theorem~2.7 is $0.1663$;
our operating value is safely above it, irrespective of the harmless
strict/non-strict convention at the printed endpoint:
\begin{equation}
   \frac16-0.1663
   =\frac{11}{30000}>0.
   \label{eq:fm-margin}
\end{equation}
Ford--Maynard also use precisely the shrunk family
$(1/2-\varepsilon,\varepsilon,\nu-2\varepsilon)$ in Theorem~2.7
\cite{FordMaynard2024}.  Thus the appearance of the third coordinate
$\nu-2\varepsilon$ is part of the published parameter geometry rather than a
new sieve hypothesis introduced here.

\subsection{A finite roughness consequence}
Write
\[
    \sigma=\nu_0-2\varepsilon.
\]
For $0\le\varepsilon\le\nu_0/100$, the Gate-0/2 finite bank proves
\begin{equation}
    \sigma\ge \frac{49}{300}.
    \label{eq:sigma-49-300}
\end{equation}
The following elementary lemma is fully independent of the deeper
Ford--Maynard analysis.

\begin{lemma}[Rough numbers have at most six prime factors]
\label{lem:omega6}
Let $n>1$ and suppose every prime divisor of $n$ is at least $n^\sigma$,
where $\sigma\ge49/300$.  Then
\[
    \Omega(n)\le6,
\]
where prime factors are counted with multiplicity.
\end{lemma}
\begin{proof}
If $k=\Omega(n)$, then
\[
   n=\prod_{j=1}^k p_j\ge n^{k\sigma}.
\]
Hence $k\sigma\le1$, so
\[
   k\le\frac1\sigma\le\frac{300}{49}<7.
\]
Since $k$ is an integer, $k\le6$.
\end{proof}
This finite statement is machine-checked in the Gate-0/2 bank.  It is useful
because the final rough $N_2$ contribution lies in a bounded-complexity class
once the required Ford--Maynard roughness identity has been supplied.

\subsection{The rough-weight input and the \texorpdfstring{$N_2$}{N2} cell sum}
The remaining source-specific rough-weight input has the form
\begin{equation}
 H(n)=(1*g)(\mathbf v(n))\,
       \mathbf 1_{\{P^-(n)\ge n^\sigma\}},
 \qquad |H(n)|\ll_{g,\nu_0}1.
 \label{eq:FM-H-rough}
\end{equation}
Ford--Maynard's Section~7 contains the corresponding rough-support mechanism;
Lemma~7.18 is used in the proof of Theorem~7.3 to restrict the support by a
least-prime-factor condition, while Theorem~8.2 combines the Section~7
machinery with a small exceptional region \cite{FordMaynard2024}.  In the
formal bank, however, the exact shifted-prime instantiation
\eqref{eq:FM-H-rough} is deliberately left uninhabited until its complete
source dictionary is certified.

The Gate-0/2 development separately records the aggregate $N_2$ cell-sum
repair: after choosing the largest prime as the final sieve variable, using a
coarse geometric mesh and aggregating the cell errors, the desired $N_2$
upper bound is reduced to the two-linear-forms upper-sieve and
prefix-volume/Mertens--PNT inputs.  The finite compiler and the
$\varepsilon$-uniformity bookkeeping are banked; the analytic source input is
not promoted to a Lean theorem.

\subsection{Gate 2 as a conditional endgame: published and proof-specific routes}
The published Ford--Maynard framework is formulated with a general Type-II
hypothesis for arbitrary divisor-bounded coefficient sequences.  We keep that
statement as the clean published sufficient wrapper and denote the required
operating-width input schematically by
\[
   \mathrm{FullFMTypeII}_{1/6}.
\]
Together with the source-exact rough $N_2$ input and the frozen Gate-0
comparison/Type-I inputs, this gives the programme-level conditional implication
\begin{equation}
\begin{aligned}
 &\mathrm{FullFMTypeII}_{1/6}
 +\mathrm{FMRoughN2}
 +\text{Gate-0 comparison inputs}\\
 &\hspace{32mm}\Longrightarrow
 \text{positive shifted-prime/twin mass at scale }x.
\end{aligned}
\label{eq:gate2-published-route}
\end{equation}
This is the published-route conditional wrapper; the analytic antecedents are
not proved by the present manuscript.

A proof audit suggests a potentially weaker route.  In the proof of
Ford--Maynard Proposition~7.22, after the finite factor decomposition the
left-hand side of their equation~(7.23) is written as a grouped bilinear sum
\[
    \sum_{n=n'n''\asymp x}Y_1(n')Y_2(n'')w_{n'n''},
\]
and the published Type-II bound is then invoked on this grouped object.  Thus
a direct theorem controlling the complete family of structured sums actually
created by the proof would be sufficient for that invocation.  We denote the
corresponding \emph{proof-specific candidate interface} by
\[
   \boxed{\texttt{FM-SIEVEGEN-TYPEII45}}.
\]
This is a weaker-hypothesis reanalysis of the proof, not a Ford--Maynard theorem
statement and not yet a completed theorem of this manuscript.

A further proposed narrowing, labelled
\[
   \boxed{\texttt{FM-PERRON-GENERATED-TYPEII45}},
\]
tracks the pre-supremum coefficients to a generated grammar consisting of
Möbius factors, box/support cutoffs, Mellin/Perron twists,
least/greatest-prime-factor phases, perfect-power pullbacks and bounded-depth
convolutions.  This generated grammar is a research target pending a complete
grammar/source certificate; it is not asserted proved.

Accordingly, a second conditional route may be written only after introducing
an explicit \emph{replacement certificate} $\mathfrak C_{\mathrm{FM}}$ which
verifies that every Type-II use in the rerun of the relevant Ford--Maynard
lower-bound proof is covered by the structured theorem.  Conditional on such a
certificate,
\begin{equation}
\begin{aligned}
 &\mathfrak C_{\mathrm{FM}}
 +\mathrm{FMRoughN2}
 +\text{Gate-0 comparison inputs}\\
 &\hspace{32mm}\Longrightarrow
 \text{positive shifted-prime/twin mass at scale }x.
\end{aligned}
\label{eq:gate2-structured-route}
\end{equation}
Neither $\mathfrak C_{\mathrm{FM}}$ nor the Perron-generated theorem is
constructed in the current paper.

\begin{theorem}[Conditional endgame, publication-safe form]
\label{thm:conditional-endgame}
Assume the frozen Gate-0 comparison/Type-I inputs and the source-exact rough
$N_2$ input.  If either
\begin{enumerate}[label=(\alph*)]
  \item the published universal Ford--Maynard Type-II hypothesis holds on the
        interval $[\varepsilon,1/6-\varepsilon]$, or
  \item a separately certified proof-specific replacement package
        $\mathfrak C_{\mathrm{FM}}$ covers every Type-II invocation required in
        the rerun of the downstream proof,
\end{enumerate}
then the corresponding conditional Gate-2 compiler yields positive
shifted-prime/twin mass at scale $x$.  If the chosen package holds at
arbitrarily large scales, the finite positivity implication yields infinitely
many distinct twin-prime pairs.
\end{theorem}
\begin{proof}
This is a conditional logical statement.  In route~(a), the published
Ford--Maynard Type-II antecedent supplies the required bilinear estimates.  In
route~(b), the replacement certificate is, by definition, required to verify
that the structured theorem covers every Type-II invocation used by the
rerun.  The rough-$N_2$ and comparison inputs supply the remaining analytic
antecedents.  The finite Gate-2 compiler then produces positive mass at the
scale.  Repeating on arbitrarily large scales yields distinct finite-support
witnesses.  The present manuscript proves neither route's missing analytic
antecedents.
\end{proof}

\subsection{Where Gate 1A and Gate 1B land in the downstream Type-II input}
The published universal Type-II wrapper remains sufficient and should not be
removed from the dependency graph.  What has changed is the programme's
candidate minimal interface.  Gate~1A supplies a conditional all-$m$ analytic
engine; the late Gate-1B architecture aims to supply the complementary
source-coherent pair-modulus multiplier estimates.  A source-exhaustive
compiler must then show that the actual generated coefficients required by the
Ford--Maynard rerun lie in the chosen structured class.

The honest dependency graph therefore has two possible Type-II entrances:
\[
\begin{gathered}
 \text{Gate 0 inputs}\qquad
 \text{Gate 1A engine}\qquad
 \text{Gate 1B pair-modulus frontier}\\
 \searrow\qquad\downarrow\qquad\swarrow\\[-1mm]
 \begin{cases}
  \mathrm{FullFMTypeII}_{1/6}, & \text{published sufficient wrapper},\\
  \mathfrak C_{\mathrm{FM}}, & \text{proof-specific replacement, if certified},
 \end{cases}\\
 \downarrow\\[-1mm]
 \text{Gate 2 Ford--Maynard endgame}\\
 \downarrow\\[-1mm]
 \text{positive twin mass at a scale}.
\end{gathered}
\]
At present, the moving source-multiplier estimate, the generated-grammar
certificate, the literal source pins and source-exhaustive reassembly remain
open.  Thus neither entrance is presently activated by the manuscript.

\subsection{Publication-facing Gate-0/2 status}
For clarity, the downstream status is summarized as follows:
\begin{center}
\begin{tabular}{p{0.43\textwidth}p{0.45\textwidth}}
\toprule
Object & Current status\\
\midrule
Gate-0 deterministic Type-I compiler & Lean-banked finite compiler\\
Gate-0 deep analytic inputs & research pass / external-uninhabited\\
$1/6>0.1663$ and margin $11/30000$ & Lean-banked arithmetic\\
$\sigma\ge49/300$ & Lean-banked arithmetic\\
$\Omega(n)\le6$ for $\sigma$-rough $n$ & Lean-banked theorem\\
Ford--Maynard rough-weight source dictionary & external/source pin; not a Lean inhabitant\\
$N_2$ cell-sum and $\varepsilon$-uniform splice & conditional research pass; Lean-uninhabited\\
Full FM Type II at width $1/6$ & OPEN / uninhabited\\
Gate 2 & CLOSED ONLY CONDITIONALLY on the displayed antecedents\\
Full Type-II reassembly & OPEN\\
Twin-prime infinitude & NOT PROVED\\
\bottomrule
\end{tabular}
\end{center}

\begin{status}[Full-programme closure status]
The role of this section is to make the strongest currently defensible
conditional closure explicit.  It does not upgrade any external,
source-inspected or programme-generated interface to a theorem.  Gate~1B is
not closed.  A subsequent source audit refines the earlier pair-modulus
frontier into a same-$q$ cross-character Gram, a cross-$q$ source-spread
certificate, and a source-linked shifted-quotient child; none is promoted to a
theorem here.  Downstream, universal Ford--Maynard Type II remains a published
sufficient wrapper, while \texttt{FM-SIEVEGEN-TYPEII45} and
\texttt{FM-PERRON-GENERATED-TYPEII45} are progressively narrower
proof-specific research targets.  The literal Gate-output/Ford-source adapter,
the generated grammar certificate, source pins and source-exhaustive
reassembly are not yet certified, and twin-prime infinitude is not proved.
\end{status}


\section{Late update: refined Gate-1B frontier and Ford--Maynard landing strip}
\label{app:late-update}

This late update records source audits completed after the main body of the
draft.  It refines the dependency graph but does not alter the statement or
proof status of the conditional Gate-1A theorem.  In particular, none of the
programme labels introduced below is promoted to a theorem merely by being
recorded here.

\subsection{Refined Gate-1B source split}
The earlier pair-modulus formulation remains a useful umbrella, but a later
source audit suggests a finer decomposition.

For the QK56 covariance parent, the cross-modulus branch $q_1\ne q_2$ lies in a
product-modulus regime in which the relevant multiplicative product ranges
are shorter than the product modulus.  In that no-wrap regime, weighted
product-energy arguments reduce the remaining issue to a source-level
spread/multiplicity statement for the push-forward multiplier.  We record the
temporary programme label
\[
   \boxed{\texttt{CROSSQ-THETA-SPREAD45}}.
\]
The subsequent internal audit recorded in
\cref{app:late-update-II} marks this source-spread step as an internal
finite/source pass, subject to final literal transcription.  It remains neither
a published theorem nor a Lean-closed analytic result.

The same-modulus branch $q_1=q_2=q$, by contrast, retains genuine modular
wraparound.  After character expansion the surviving object has schematic
shape
\[
 \sum_{q\asymp Q}|\beta(q)|^2
 \sum_{\chi_1,\chi_2\bmod q}
 \widehat R_q(\chi_1)\,
 \overline{\widehat R_q(\chi_2)}\,
 \mathcal W_q(\chi_1,\chi_2).
\]
The diagonal $\chi_1=\chi_2$ does not exhaust this Gram, and fixed-multiplier
Kloosterman estimates do not remove the cross-character terms.  The subsequent internal diagonal/off-diagonal audit recorded in
\cref{app:late-update-II} gives this child an internal analytic pass, with the
literal QK56/S1/S2 source dictionary still load-bearing:
\[
   \boxed{\texttt{SAMEQ-QK56-CROSSCHAR-GRAM45: internal analytic pass}}.
\]
This remains a programme claim subject to source transcription and hostile
audit, not an externally published theorem.

The balanced shifted-quotient branch likewise carries source-linked four-cycle
multipliers rather than four independent arbitrary multipliers.  Its first
required action is now the literal source recovery
\[
   \boxed{\texttt{SHIFT-TTSTAR-LITERAL-SOURCE-PIN}},
\]
after which the narrower research child
\[
   \boxed{\texttt{SHIFT-SOURCE-LINKED-CHAR45}}
\]
is the first major analytic candidate.  Published
fixed-modulus Kloosterman results remain backend-capacity inputs only:
Blomer--Pascadi give a $c^{-1/32}$ saving in the square-root critical range
\cite{BlomerPascadi2026}, while Pascadi's non-abelian amplification gives
$c^{-1/12}$ for balanced products of two comparable primes
\cite{Pascadi2025}.  Neither published statement controls the above moving
source-linked family automatically.

Thus the refined Gate-1B research ledger is
\[
\begin{array}{ll}
 \texttt{SAMEQ-QK56-CROSSCHAR-GRAM45}
   & \text{internal analytic pass; source pin pending},\\
 \texttt{CROSSQ-THETA-SPREAD45}
   & \text{internal finite/source pass; transcription pending},\\
 \texttt{SHIFT-TTSTAR-LITERAL-SOURCE-PIN}
   & \text{source first action open},\\
 \texttt{SHIFT-SOURCE-LINKED-CHAR45}
   & \text{first major analytic open},\\
 S1/S2\text{ literal normalisation}
   & \text{narrow source pins},\\
 \text{degeneracy/collision/packet census}
   & \text{finite routing/reassembly}.
\end{array}
\]
This supersedes only the research-level description of the common Gate-1B
frontier.  Gate~1B remains open.

\subsection{Published Ford--Maynard five-collar geometry}
One downstream geometric source issue can be removed from the open ledger.
For
\[
   P_\varepsilon=
   \left(\frac12-\varepsilon,\,
         \varepsilon,\,
         \nu-2\varepsilon\right),
\]
Ford--Maynard's proof of Theorem~8.3 shows that every vector
$x\in H_2$ has a component of one of the five forms
\begin{equation}
 \frac12+u,\qquad
 \nu+u,\qquad
 2\nu+u,\qquad
 1-\nu+u,\qquad
 1-2\nu+u,
 \qquad |u|\le 6\varepsilon.
 \label{eq:FM-five-collars}
\end{equation}
This is a published source statement \cite{FordMaynard2024}.  Moreover,
Ford--Maynard Definition~5.1 defines
\[
   \mathbf v(n)=
   \left(
     \frac{\log p_1}{\log n},\ldots,
     \frac{\log p_k}{\log n}
   \right),
   \qquad n=p_1\cdots p_k,
\]
so the coordinates in \eqref{eq:FM-five-collars} correspond to individual
prime factors.  At $\nu_0=1/6$ the collar centres are
\[
   \frac12,\qquad \frac16,\qquad \frac13,\qquad
   \frac56,\qquad \frac23.
\]
We therefore record the programme status
\[
   \boxed{\texttt{FM-H2-FIVE-COLLARS45: published-source pass}.}
\]

\subsection{A source-specific repair of the rough \texorpdfstring{$N_2$}{N2} sector}
The common normalization $a_n,b_n\mapsto a_n/\log x,b_n/\log x$ does not by
itself place the shifted-prime comparison pair into the bounded Ford--Maynard
class: it also reduces the prime mass by an additional factor of $\log x$.
Accordingly, that normalization is not used as the $N_2$ repair.

Instead set
\[
   \nu_0=\frac16,\qquad \sigma=\nu_0-2\varepsilon.
\]
For $0<\varepsilon\le \nu_0/100$ one has
$\sigma\ge49/300$, and the finite bank proves that every $\sigma$-rough
integer has at most six prime factors counted with multiplicity.

The squareful rough subfamily is negligible.  Indeed, if $p^2\mid n\le x$
and $p\ge x^\sigma$, then
\begin{equation}
 \sum_{x^\sigma\le p\le x^{1/2}}
 \left(\frac{x}{p^2}+1\right)\log x
 \ll
 x^{1-\sigma}\log x+x^{1/2}\log x
 =o\!\left(\frac{x}{\log x}\right).
 \label{eq:N2-squareful}
\end{equation}
For the squarefree sector, the five-collar geometry permits one to choose
canonically a collar prime $q$ and then the largest remaining prime $r$, writing
$n=qdr$.  The present research reduction controls the shifted-prime part by a
uniform upper-bound sieve for the two linear forms
\[
     r,\qquad qdr+2.
\]
On the rough sector the associated local factor is uniformly bounded, while
the harmonic sum over the narrow prime collar contributes
$O_{\nu_0}(\varepsilon+1/\log x)$.  This yields the research-level estimate
\[
   \sum_{n\in N_2} a_n|H(n)|
   \ll
   \left(\varepsilon+o(1)\right)\frac{x}{\log x},
\]
conditional on the standard uniform two-linear-forms upper-sieve input in the
stated rough range.  The comparison side is compatible with the
Ford--Maynard rough-number count and the same five-collar volume estimate.
We record this as
\[
 \boxed{\texttt{FM-N2-SHIFTED-PRIME-COLLAR-SIEVE45}}
\]
with research status ``pass modulo a standard uniform two-linear-forms upper
sieve''.  This remains external/uninhabited in Lean until the literal shifted-prime
source dictionary and analytic upper-sieve input are instantiated.

\subsection{The proof-specific Ford--Maynard landing point}
The published Ford--Maynard Type-II hypothesis remains the clean sufficient
wrapper.  However, Proposition~7.22 provides a more economical
\emph{proof-specific} landing point.  After the finite factor decomposition,
Ford--Maynard reduce the relevant estimate to their equation~(7.23), namely a
finite-depth product sum with one selected product constrained to the Type-II
interval.  They then group the factors into
\[
   Y_1(n')
   =
   \sum_{n'=\prod_{e\in E}n_e}
      \prod_{e\in E}\mathfrak x_e(n_e),
   \qquad
   Y_2(n'')
   =
   \sum_{n''=\prod_{e\notin E}n_e}
      \prod_{e\notin E}\mathfrak x_e(n_e),
\]
and invoke the universal Type-II bound only after reaching
\[
   \sum_{n=n'n''\asymp x}
      Y_1(n')Y_2(n'')w_{n'n''}.
\]
This is a direct reading of the proof of Proposition~7.22
\cite{FordMaynard2024}.  Consequently, a theorem controlling \emph{all actual
generated instances} of Ford--Maynard equation~(7.23) is sufficient for this
specific invocation:
\[
 \boxed{
 \text{actual generated Ford--Maynard (7.23)}
 \Longrightarrow
 \text{Proposition 7.22}
 \Longrightarrow
 \text{Proposition 7.19}.}
\]
This is a proof audit, not a new theorem of Ford--Maynard and not a claim that
the antecedent has been proved here.

The factor depth is fixed by the sieve parameters.  In the proof,
\[
   N=k\ell+r+s,
\]
with $k,\ell,r,s$ bounded in terms of the fixed sieve data; thus the earlier
concern that the generated landing point necessarily has unbounded convolution
depth is not an independent obstruction.

We retain the temporary label
\[
   \boxed{\texttt{FM-PERRON-GENERATED-TYPEII45}}
\]
with status:
\[
   \boxed{\text{proof-specific sufficient target; generated grammar/source certificate open}.}
\]

\subsection{The remaining Gate-output/Ford-source adapter}
The present Gate-1A theorem controls a canonical common physical weight source,
and the Gate-1B analysis has its own literal source classes.  Ford--Maynard's
factor decomposition, however, generates a finite family of one-variable
factors before equation~(7.23).  The next downstream question is therefore
first a source dictionary problem:
\[
   \omega_{\Omega,j}\longmapsto
   \begin{cases}
     \text{Gate-1A source},\\
     \text{Gate-1B source},\\
     \text{direct exceptional source}.
   \end{cases}
\]
We use the programme label
\[
   \boxed{\texttt{SOURCE-ADAPTER45}}
\]
for this literal Gate-output/Ford-source splice.  Its current status is:
\[
   \boxed{\text{primary candidate downstream interface; source-exhaustive dictionary open}.}
\]
Until the literal audit is complete, it is not known whether this is merely a
finite compiler problem or whether one generated weight family exposes a
stronger analytic child.

\subsection{Revised downstream ledger}
The late-update status is:
\[
\begin{array}{ll}
 \text{Ford }H_2\text{ five-collar geometry}
   & \text{published-source pass},\\
 \sigma\text{-rough }\Omega(n)\le6
   & \text{Lean-banked finite theorem},\\
 N_2\text{ squareful disposal}
   & \text{pass},\\
 N_2\text{ shifted-prime two-form lane}
   & \text{research pass modulo standard upper sieve},\\
 \text{Ford (7.23) proof-specific landing}
   & \text{proof-audit pass},\\
 \text{unbounded-depth grammar concern}
   & \text{retired},\\
 \texttt{SOURCE-ADAPTER45}
   & \text{primary downstream candidate interface},\\
 \texttt{SAMEQ-QK56-CROSSCHAR-GRAM45}
   & \text{internal analytic pass; source pins pending},\\
 \texttt{SHIFT-SOURCE-LINKED-CHAR45}
   & \text{first major Gate-1B analytic open},\\
 \text{Gate 1B}
   & \textbf{OPEN},\\
 \text{Full Type-II reassembly}
   & \textbf{OPEN},\\
 \text{twin-prime infinitude}
   & \textbf{NOT PROVED}.
\end{array}
\]

The revised shortest dependency graph is
\[
\begin{gathered}
 \text{Gate-0 Type-I/comparison inputs}
 +\text{Gate-1A literal output}
 +\text{Gate-1B literal output}\\
 \downarrow\\
 \boxed{\texttt{SOURCE-ADAPTER45 / source-exhaustive Ford-(7.23) compiler}}\\
 \downarrow\\
 \text{actual generated Ford--Maynard (7.23)}\\
 \downarrow\\
 \text{Proposition 7.22}\to\text{Proposition 7.19}
 +\text{source-specific }N_2\text{ repair}\\
 \downarrow\\
 \text{Ford--Maynard lower-bound rerun at }\nu=1/6\\
 \downarrow\\
 \text{positive shifted-prime/twin mass at a scale}.
\end{gathered}
\]
Every arrow after the source-exhaustive compiler remains conditional on the
literal analytic antecedents displayed above.  This diagram does not assert
the Twin Prime Conjecture.

\begin{status}[Late-update publication firewall]
The late update narrows the downstream uncertainty, but does not reduce the
programme to one remaining theorem.  In particular:
\begin{itemize}
  \item Gate~1B is not closed;
  \item \texttt{SAMEQ-QK56-CROSSCHAR-GRAM45} is an internal analytic pass
        subject to literal source pins and hostile audit, not a published or
        Lean-banked analytic theorem;
  \item \texttt{SOURCE-ADAPTER45} has not yet been classified as finite-only
        versus genuinely analytic;
  \item \texttt{FM-PERRON-GENERATED-TYPEII45} is proof-specific and the
        generated source grammar is not yet certified;
  \item the Ford--Maynard endgame remains conditional;
  \item twin-prime infinitude is not proved.
\end{itemize}
\end{status}


\section{Late update II (historical): QK56 promotion, shifted source-linked child, and source-exact Ford adapter}
\label{app:late-update-II}

This section records an intermediate late-26/27-August programme audit.  It
does not alter Theorem~\ref{thm:main} and does not assert Gate-1B closure, Full
Type-II closure, or twin-prime infinitude.  The subsequent native first-$TT^*$
derivation in \cref{app:native-ttstar} supersedes the stronger full-covariance
status wording below: only the same-$q$ cross-character child and the
cross-$q$ operator theorem survive as positive analytic children; the
same-$q$ diagonal source energy and cross-$q$ physical source spread remain
open.

\subsection{QK56 full covariance: internal promotion}
The same-modulus covariance is first split into the character diagonal and the
cross-character part.  Schematically,
\[
 \sum_{q\sim Q}|\beta(q)|^2
 \sum_{\chi_1,\chi_2\bmod q}
 \widehat R_q(\chi_1)\overline{\widehat R_q(\chi_2)}
 \mathcal W_q(\chi_1,\chi_2)
 =\mathcal G_{\rm diag}+\mathcal G_{\rm off}.
\]
The internal audit places the large $HK/q$ contribution in
$\mathcal G_{\rm diag}$ and routes that diagonal to the existing
coefficient-blind residue-energy/JQ7 bank.  The off-diagonal is identified with
a centred modular-hyperbola discrepancy and is estimated internally through a
square-root Kloosterman/Weil discrepancy argument.  The current worst signed
programme margin is
\[
       X^{-5/36+o(1)}.
\]
We therefore record
\[
 \boxed{\texttt{SAMEQ-QK56-CROSSCHAR-GRAM45: internal analytic pass}}
\]
subject to the literal S1/S2 source pin and final hostile audit.

For $q_1\ne q_2$, the latest no-wrap/source audit uses the fact that fixing the
source parameter $\Theta$ determines $B_1\bmod q_1$ and $B_2\bmod q_2$.
The remaining product/numerical multiplicity is internally within the required
$\ell^2/\ell^1$ spread budget.  We therefore record
\[
 \boxed{\texttt{CROSSQ-THETA-SPREAD45: historical candidate; source spread later reopened}},
\]
again without promoting the statement to a published theorem or a Lean-closed
analytic result.

At programme level the combined status is
\[
 \boxed{\begin{gathered}
   \texttt{QK56-FULL-COVARIANCE45}\\
   \text{historical programme promotion; superseded by Late update III}
 \end{gathered}}
\]

\subsection{The shifted child: source action before analysis}
The balanced shifted-quotient branch should not be modelled as four independent
multipliers.  Its cyclic source relation has the form
\[
    a_j=\Theta(B_j,B_{j+1})\qquad(j\bmod4),
\]
and the schematic character moment is
\[
 \sum_{B_1,\ldots,B_4} W(\mathbf B)
 \sum_h W_{\mathbf B}(h)\,
 \chi\!\left(
   \mathscr D\bigl(
      \Theta(B_1,B_2),\Theta(B_2,B_3),
      \Theta(B_3,B_4),\Theta(B_4,B_1);h
   \bigr)
 \right).
\]
The four-cycle discriminant algebra is banked in the programme, but the literal
pre-support-enlargement $TT^*$ source formula has not yet been recovered.
Consequently the first action is
\[
 \boxed{\texttt{SHIFT-TTSTAR-LITERAL-SOURCE-PIN}},
\]
which must print the literal $B_i$ ranges, the exact cyclic $\Theta$ map, the
$h$-range and weight, prefactor and normalization, physical target, all
zero/nonunit/collision/repeated-prime strata, and labelled multiplicity.
Only after this transcription is the analytic target
\[
 \boxed{\texttt{SHIFT-SOURCE-LINKED-CHAR45}}
\]
well posed.  It is the current first major analytic Gate-1B candidate.

\subsection{Technology cards for the shifted source-linked child}
Several published results are relevant only after the literal source has been
reduced to their stated geometries.

Fouvry--Shparlinski estimate weighted multiplicative-character sums of the
$2\times2$ determinant $ad-bc$ over four interval variables, with power saving
for $N\ge p^{1/8+\varepsilon}$ \cite{FouvryShparlinski2022}.  This is a serious
candidate if the cyclic discriminant is reduced source-exactly to determinant
geometry.

Shparlinski treats multivariate exponential and multiplicative-character sums
with monomials, allowing bounded one-variable weights and both positive and
negative monomial exponents \cite{Shparlinski2013}.  This is relevant if the
cyclic $\Theta$ dictionary reduces to a monomial or reciprocal-monomial kernel.
Bourgain--Garaev's reciprocal sumset and multilinear Kloosterman technology is
a secondary candidate once an actual reciprocal-product dictionary has been
manufactured \cite{BourgainGaraev2012}.

The very general trace-function bilinear theorem of
Fouvry--Kowalski--Michel--Sawin requires substantial structural hypotheses on
the geometric monodromy group (for instance a simple connected identity
component, or finite quasisimple monodromy in the displayed simplified
theorem) \cite{FouvryKowalskiMichelSawin2025}.  A rank-one Kummer/Legendre
character does not satisfy those displayed hypotheses merely by being a trace
function, so that result is not cited as a black box for the present child.
Blomer--Pascadi remains a fixed-Kloosterman backend card and likewise does not
resolve the source-linked discriminant family by itself
\cite{BlomerPascadi2026}.

\subsection{Minimal current Gate-1B ledger}
The late-update Gate-1B ledger is therefore
\begin{center}
\small
\begin{tabular}{p{0.45\linewidth}p{0.43\linewidth}}
\texttt{QK56-FULL-COVARIANCE45}
  & historical promotion; full source-exact closure reopened\\
\texttt{SHIFT-TTSTAR-LITERAL-SOURCE-PIN}
  & source open; first action\\
\texttt{SHIFT-SOURCE-LINKED-CHAR45}
  & first major analytic open\\
S1/S2 literal normalization
  & narrow source pins\\
$D=0$/shared-prime/nonunit/collision census
  & finite closure/routing\\
fixed/switched reassembly
  & finite compiler
\end{tabular}
\end{center}
Thus Gate~1B is still not closed.  The later native first-$TT^\ast$ source
derivation shows that QK56 itself still contains two source-exact quantitative
gaps---the same-$q$ character diagonal and cross-$q$ CRT source spread---before
the shifted-source branch and finite reassembly can be assessed exhaustively.

\subsection{Ford-generated one-variable archaeology and the Gate-output source dictionary}
This source question is separate from Gate-1B analysis.  The latest internal
Ford-side archaeology records the following programme passes:
\[
\begin{array}{ll}
 \texttt{FM-GENERATED-ONEVAR-ARCHAEOLOGY45} & \text{pass},\\
 \texttt{FM-FINITE-NUCLEAR-COMPILER45} & \text{pass},\\
 \texttt{FM-OUTER-COEFFICIENT-ABSORPTION45} & \text{pass},\\
 \texttt{WEIGHT-DEPENDENCE-FINITE-COMPILER45} & \text{pass}.
\end{array}
\]
These are programme-level source/compiler audits, not published theorems.
The first unresolved interface is the source-exact weighted high-complexity
packet dictionary.  For each actual packet it must provide the operator,
ranges, coefficient, multiplicity, weight law, centering, gcd/collision
restrictions, original prefactor, normalized target and route.

In particular the actual Gate weight law is not to be guessed: the relevant
packet must be identified as a common $W_D$, a finite-template
$W_{D,e}$, or a genuinely edge-dependent $W_{D,e}$.  We therefore retain
\[
   \boxed{\texttt{FM-TO-GATE-COORDINATE-CENSUS45: source-blocked}}
\]
rather than calling it a new analytic theorem.  Once this literal source
census exists, the finite coordinate/reassembly compiler is already available
at programme level.

\subsection{Ford--Maynard scope}
Ford--Maynard's published framework remains the general sufficient wrapper for
non-negative sequences satisfying its Type-I/Type-II hypotheses
\cite{FordMaynard2024}.  The programme's generated-one-variable and generated
(7.23) routes are proof-specific weakenings under investigation; we do not
attribute them to Ford--Maynard as theorem statements.  In particular, the
latest source/compiler audits do not change the final publication firewall:
\[
 \boxed{\begin{gathered}
   \text{Gate 1B open; downstream source dictionary open;}\\
   \text{Ford--Maynard endgame conditional; twin primes not proved.}
 \end{gathered}}
\]

\begin{status}[Draft 10 late-update firewall]
The same-$q$ cross-character promotion remains an internal programme result;
the cross-$q$ operator theorem is likewise internal and does not provide the
physical source-spread estimate.  The stronger QK56 full-covariance status in
this historical section is superseded by \cref{app:native-ttstar}.  None of
these statements is a kernel-checked Lean analytic theorem or an externally
published Gate-1B theorem.  Likewise, \texttt{SHIFT-SOURCE-LINKED-CHAR45}, the
physical high-complexity packet dictionary, the FM-to-Gate coordinate census
and the Ford lower-bound rerun must not be promoted without their actual
proof/source objects.
\end{status}



\section{Late update III (partial child analysis): native first-\texorpdfstring{$TT^\ast$}{TT*} Gate-1B reduction}
\label{app:native-ttstar}

This section supersedes the stronger QK56 full-covariance status language in
\cref{app:late-update-II}.  It is a source-explicit programme update, not a
claim of Gate-1B closure.  The decisive correction is that the first
$TT^\ast$ covariance has three different modulus geometries.  The two
positive analytic children below remain valid on their stated proper-smooth
sectors.  Late update~IV supersedes this section as the controlling
high-conductor architecture: once the complement M\"obius sign is restored,
same-$q$ and common-conductor cross-$q$ can no longer be closed sequentially.

\subsection{Scope and current status}

This note is deliberately not a claim of Gate--1B closure.  It isolates the
part of the remaining analysis that can now be written in source-exact
coordinates and separates it from source-definition and reassembly tasks.

The controlling internal status is:
\[
\begin{array}{ll}
\text{same-}q\text{ smooth cross-character discrepancy}
  & \text{analytic pass on the proper smooth sector},\\
\text{coprime cross-}q\text{ no-wrap moving-multiplier }L^2
  & \text{proved at the kernel/operator level},\\
\text{same-}q\text{ character diagonal}
  & \text{source-energy open},\\
\text{cross-}q\text{ CRT source spread}
  & \text{source-norm open},\\
\text{Full-Nine/Ford--Maynard reassembly}
  & \text{source-exhaustive compiler open},\\
\text{shifted four-cycle literal }TT^\ast\text{ source}
  & \text{source definition open}.
\end{array}
\]

Published Kloosterman and determinant-character technology is useful only
after the literal source dictionary is fixed.  In particular, fixed-modulus
bilinear Kloosterman estimates do not by themselves supply a moving source
multiplier family theorem.

\subsection{Native order-five source and two-slot completion}

Let
\[
 Y=X^{1/9},\qquad q\asymp Q=X^\omega,\qquad
 \frac{13}{18}\le \omega<\frac89,\qquad qR\asymp X=Y^9.
\]
In the critical order-five cell, write the product of five defect factors as
$C_5\asymp Y^5$ and retain four model variables
$x_1,x_2,x_3,x_4\asymp Y$.  Put
\[
   B=C_5x_2x_3\asymp Y^7.
\]
The native shifted relation is
\begin{equation}
   \boxed{Bx_1x_4-q\ell=-2.}
   \label{eq:native-shell}
\end{equation}
This is the determinant shell generated by the $n+2$ source; it is not a
replacement by a $2n\pm1$ model.

For fixed $(B,q)$ the two physical Poisson variables are $x_1,x_4$.  With
smooth source factors $f_1,f_4,W$, set
\[
 F_{B,q}(x_1,x_4)
 =
 f_1(x_1)f_4(x_4)
 W\!\left(\frac{Bx_1x_4+2}{qR}\right)
 \1_{Bx_1x_4\equiv-2\pmod q}.
\]
Two-dimensional completion gives
\begin{equation}
 \mathcal T_{B,q}
 =
 \frac{Y^2}{q^2}
 \sum_{h,k\in\Z}
   \Psi_{B,q}(h,k)
   S(h,-2k\overline B;q).
 \label{eq:two-poisson}
\end{equation}
On the proper unit sector $(kB,q)=1$, Kloosterman scaling yields
\[
  S(h,-2k\overline B;q)=S(hk\overline B,-2;q),
\]
so the actual one-copy multiplier is
\begin{equation}
   \boxed{a=hk\overline B\pmod q.}
   \label{eq:one-copy-multiplier}
\end{equation}
The effective ranges are
\[
   |h|,|k|\ll q/Y,
\]
and the normalized smooth kernel has uniformly bounded derivatives on a
fixed compact box.  Hence the smooth source can be separated into
polylogarithmically controlled nuclear packets.  This separation is a
source-normalization device, not an independent arithmetic saving.

\subsection{The first \texorpdfstring{$TT^\ast$}{TT*} sector law}

For one nuclear packet write
\[
 \mathcal F_q(h,k)
 =
 \sum_B c_{q,B}(h,k)S(hk\overline B,-2;q).
\]
Opening the covariance produces
\[
 \sum_{h,k}\left|\sum_q\mathcal F_q(h,k)\right|^2
 =
 \sum_{q_1,q_2,B_1,B_2,h,k}
 c_{q_1,B_1}(h,k)\overline{c_{q_2,B_2}(h,k)}
 \mathcal K_{12}(hk),
\]
where
\[
 \mathcal K_{12}(t)
 =
 S(t\overline{B_1},-2;q_1)
 \overline{S(t\overline{B_2},-2;q_2)}.
\]

There are three genuinely different arithmetic sectors.

\subsubsection{Same modulus}

If $q_1=q_2=q$, the product is most naturally diagonalized by
multiplicative characters.  For $(mn,q)=1$,
\begin{equation}
 S(m,n;q)
 =
 \frac1{\phi(q)}
 \sum_{\chi\bmod q}
   \tau_q(\chi)^2\overline{\chi}(mn).
 \label{eq:kloos-char}
\end{equation}
Thus the same-$q$ covariance is a character Gram rather than a single
moving Kloosterman sum.  Its character diagonal and cross-character pieces
must be treated separately.

\subsubsection{Coprime cross-modulus}

If $q_1\ne q_2$ and $(q_1,q_2)=1$, CRT multiplicativity gives
\begin{equation}
 \boxed{\mathcal K_{12}(t)=S(t,\Theta_{12};q_1q_2),}
 \label{eq:crossq-kloos}
\end{equation}
where $\Theta_{12}\pmod{q_1q_2}$ is uniquely determined by
\[
 \Theta_{12}\equiv-2q_2^2\overline{B_1}\pmod{q_1},
 \qquad
 \Theta_{12}\equiv-2q_1^2\overline{B_2}\pmod{q_2}.
\]
This is the actual first-$TT^\ast$ moving multiplier.

\subsubsection{Shared divisor}

When $(q_1,q_2)>1$, the product modulus cannot be treated by
\cref{eq:crossq-kloos} without first extracting the common conductor.
This branch belongs to the existing GCD/shared-prime router.

\begin{proposition}[First-$TT^\ast$ sector law]
The native QK56 covariance separates into:
\[
\begin{aligned}
  \text{same-}q&=\text{character Gram},\\
  (q_1,q_2)=1&=\text{CRT Kloosterman},\\
  (q_1,q_2)>1&=\text{shared-conductor router}.
\end{aligned}
\]
In particular, the first covariance is not a universal four-independent-
multiplier character moment.
\end{proposition}

\subsection{Same-\texorpdfstring{$q$}{q} cross-character discrepancy}

The cross-character part of the same-$q$ Gram admits a source-specific
modular-hyperbola treatment.

Let $G_q=(\Z/q\Z)^\times$.  For a pair of separated smooth frequency
weights $u,v$ define
\[
  g_q(r)
  =
  \sum_{\substack{h,k\in G_q\\hk=r}}u(h)v(k)
\]
and
\[
  m_q
  =
  \frac1{\phi(q)}
  \left(\sum_hu(h)\right)
  \left(\sum_kv(k)\right).
\]
Set
\[
   \Delta_q(u,v)=\max_{r\in G_q}|g_q(r)-m_q|.
\]

\subsubsection{Prime modulus}

For prime $q$, additive completion gives the exact identity
\[
 g_q(r)
 =
 \frac1{q^2}
 \sum_{a,b\bmod q}
   \widehat u(a)\widehat v(b)S(a,br;q).
\]
After removing the zero frequencies one obtains
\[
 g_q(r)-\frac{UV}{q-1}
 =
 \frac1{q^2}
 \sum_{\substack{a,b\bmod q\\a,b\ne0}}
   \widehat u(a)\widehat v(b)S(a,br;q)
 -
 \frac{UV}{q^2(q-1)},
\]
where $U=\sum u$ and $V=\sum v$.

Define the discrete bounded-variation seminorm
\[
  \mathcal B(u)
  =
  \norm{u}_\infty+\sum_x|u(x+1)-u(x)|.
\]
Summation by parts gives
\[
 \sum_{a\ne0}|\widehat u(a)|
 \ll
 q\,\mathcal B(u)\log(2q),
\]
and similarly for $v$.  The Weil bound then yields
\begin{equation}
 \Delta_q(u,v)
 \ll
 q^{1/2}\mathcal B(u)\mathcal B(v)\log^2(2q)
 +
 O\!\left(\frac{|UV|}{q^3}\right).
 \label{eq:prime-discrepancy}
\end{equation}

\subsubsection{Composite modulus}

For general proper-unit $q$, the modular inverse graph has Fourier
coefficients $S(a,br;q)$.  Combining the standard composite Kloosterman
bound with a two-dimensional Erd\H{o}s--Tur\'an--Koksma inequality and
partial summation gives the source-specific estimate
\begin{equation}
   \boxed{
   \Delta_q(u,v)
   \ll
   q^{1/2+o(1)}
   \mathcal B(u)\mathcal B(v).
   }
   \label{eq:composite-discrepancy}
\end{equation}
For the actual nuclear smooth weights,
$\mathcal B(u),\mathcal B(v)\ll(\log X)^{O(1)}$.

\begin{status}
Equation \eqref{eq:composite-discrepancy} is an internal analytic
derivation for the proper smooth source.  A formal proof assistant
development should separate the finite Fourier identities from the external
Kloosterman/ETK discrepancy input rather than encode the latter as an axiom.
\end{status}

\subsubsection{Fixed-power margin}

For a nonprincipal ratio character $\eta=\psi\overline\chi\ne1$, the
constant term $m_q$ cancels, and the cross-character convolution operator
has norm
\[
   \norm{T_{\mathrm{off}}}_{\mathrm{op}}
   =
   \phi(q)\Delta_q(u,v)
   \ll
   q^{3/2}(\log X)^{O(1)}.
\]
The natural uncentered $(h,k)$ mass is
\[
   HK\asymp(Q/Y)^2.
\]
Hence the ratio is
\[
   \frac{q^{3/2}}{HK}
   \asymp
   \frac{Y^2}{Q^{1/2}}
   =
   X^{2/9-\omega/2}.
\]
At $\omega\ge13/18$,
\[
   2/9-\omega/2\le-5/36.
\]
Thus
\begin{equation}
 \boxed{
 \text{same-}q\text{ cross-character contribution}
 \ll
 X^{-5/36+o(1)}
 \times
 \text{natural source energy}.
 }
 \label{eq:sameq-margin}
\end{equation}
This closes the proper smooth cross-character child, but not the character
diagonal.

\subsection{Cross-\texorpdfstring{$q$}{q} no-wrap moving-multiplier theorem}

Let
\[
 c=q_1q_2,\qquad (q_1,q_2)=1,
\]
and for $D\in G_c$ define
\[
 T_c(D)
 =
 \sum_{\substack{|h|\le H\\|k|\le K}}
 \alpha_h\beta_kS(Dh,k;c).
\]
Assume
\begin{equation}
   2HK<c.
   \label{eq:no-wrap}
\end{equation}
In the actual regime
\[
   H,K\asymp Q/Y,\qquad c\asymp Q^2,
\]
so \eqref{eq:no-wrap} holds for large $X$.

Character diagonalization gives
\[
 T_c(D)
 =
 \frac1{\phi(c)}
 \sum_{\chi\bmod c}
 \tau_c(\chi)^2\overline{\chi}(D)A_\chi B_\chi,
\]
with
\[
  A_\chi=\sum_h\alpha_h\overline{\chi}(h),
  \qquad
  B_\chi=\sum_k\beta_k\overline{\chi}(k).
\]
Orthogonality in $D$ yields
\[
 \sum_{D\in G_c}|T_c(D)|^2
 =
 \frac1{\phi(c)}
 \sum_{\chi\bmod c}
 |\tau_c(\chi)|^4|A_\chi|^2|B_\chi|^2.
\]
The character fourth moment expands into
\[
 \phi(c)
 \sum_{h_1k_1\equiv h_2k_2\pmod c}
 \alpha_{h_1}\beta_{k_1}
 \overline{\alpha_{h_2}\beta_{k_2}}.
\]
By \eqref{eq:no-wrap}, the congruence lifts to the integer equality
\[
   h_1k_1=h_2k_2.
\]
The remaining multiplicative-convolution energy is bounded by divisor
multiplicity:
\[
 \sum_n\left|\sum_{hk=n}\alpha_h\beta_k\right|^2
 \ll
 c^{o(1)}
 \norm{\alpha}_2^2\norm{\beta}_2^2.
\]
Together with the standard Gauss-sum bound,
\begin{equation}
 \boxed{
 \sum_{D\in G_c}|T_c(D)|^2
 \ll
 c^{2+o(1)}
 \norm{\alpha}_2^2\norm{\beta}_2^2.
 }
 \label{eq:crossq-l2}
\end{equation}
Consequently, for any moving-source coefficient $A(D)$,
\begin{equation}
 \boxed{
 \left|\sum_{D\in G_c}A(D)T_c(D)\right|
 \ll
 c^{1+o(1)}
 \norm{A}_2\norm{\alpha}_2\norm{\beta}_2.
 }
 \label{eq:crossq-moving}
\end{equation}

\begin{remark}
The analytic kernel in \eqref{eq:crossq-moving} is therefore controlled
without a fixed-multiplier-to-family promotion.  The remaining issue is the
physical CRT push-forward norm of $A(D)$.
\end{remark}

\subsection{A route that can be retired}

A direct completion in the whole products $C,a$ formally gives
\[
 \sum_{C,a}F(C)G(a)\1_{Ca\equiv-2\pmod q}
 =
 \frac1{q^2}
 \sum_{h,k\bmod q}
 \widehat F_q(h)\widehat G_q(k)S(h,-2k;q).
\]
However, $F$ and $G$ are sparse divisor convolutions, not interval-smooth
weights, so their Fourier transforms need not localize to short dual
intervals.

Even under the counterfactual assumption of such localization, at the
endpoint
\[
 q=X^{13/18},\qquad C=X^{5/9},\qquad a=X^{4/9},
\]
one has
\[
 q/C=q^{3/13},\qquad q/a=q^{5/13}.
\]
The most favorable fixed-multiplier contraction supplied by the relevant
bilinear Kloosterman range is still far smaller than the physical
normalization requirement.  Thus the whole-product route is nonclosing even
before addressing the moving-source family.  The slot-preserving
$(x_1,x_4)$ completion and the same-/cross-$q$ sector split are the correct
architecture.

\subsection{The exact remaining closure system}

The preceding analysis removes the need for a mysterious generic
moving-multiplier theorem.  The remaining obligations are source-exact.

\subsubsection{Same-\texorpdfstring{$q$}{q} character diagonal}

For $\eta=\psi\overline\chi=1$, the centred modular-hyperbola discrepancy
does not provide cancellation.  The required estimate has the schematic
form
\[
 \sum_{q\sim Q}
 \frac{|\beta(q)|^2Y^4}{q^4\phi(q)^2}
 HK
 \sum_{\chi\bmod q}
 |\tau_q(\chi)|^4|A_{\nu,q}(\chi)|^2
 \ll_A
 \frac{X^2Y^2}{Q^2}(\log X)^{-2A}.
\]
This must be proved using the literal seven-box source energy, the exact
normalization of $\beta(q)$, and the existing residue/JQ7 router.  It is not
a consequence of the cross-character estimate.

\subsubsection{Cross-\texorpdfstring{$q$}{q} physical source spread}

The analytic operator is controlled by \eqref{eq:crossq-moving}.  What
remains is to prove for the physical CRT push-forward
\[
 (B_1,B_2)\longmapsto \Theta_{12}
\]
a source norm of the form
\begin{equation}
  \boxed{
  \norm{A_{q_1,q_2}}_2
  \le
  Y^{-1+o(1)}
  \norm{A_{q_1,q_2}}_1,
  }
  \label{eq:source-spread}
\end{equation}
or any stronger equivalent estimate after restoring the full prefactor.
This statement must include all labelled multiplicities and shared-prime
strata.

\subsubsection{Source-exhaustive Full-Nine/Ford reassembly}

Closing the QK56 parent does not by itself prove that all balanced Full-Nine
packets have been covered.  The source compiler must establish an exhaustive
packet identity and then map the actual generated Ford--Maynard grouped
Type-II instances into Gate-1A, Gate-1B, or explicitly controlled
exceptional packets.

A parallel source-definition obligation remains for the shifted four-cycle
branch.  Existing archives contain the four-cycle matrix/discriminant
algebra, but not the literal pre-support-enlargement $TT^\ast$ physical
source.  Before a source-linked character theorem is well posed one must
recover:
\[
\begin{aligned}
 &B_1,\ldots,B_4,\qquad \Theta(B_i,B_{i+1}),\\
 &h_i\text{-ranges and weights},\qquad \text{prefactor/target},\\
 &\text{degenerate strata},\qquad \text{labelled multiplicity}.
\end{aligned}
\]

\subsection{Minimal honest dependency graph}

The current Gate--1B closure path is therefore:
\[
\begin{array}{c}
\text{same-}q\text{ diagonal source energy}\\
+\text{cross-}q\text{ CRT source spread}\\
+\text{literal S1/S2/source normalizations}\\
+\text{finite shared-prime/nonunit/collision routing}\\[1mm]
\downarrow\\
\text{QK56 full covariance closed}\\[2mm]
+\text{actual switched nine-coordinate dictionary}\\
+\text{literal shifted }TT^\ast\text{ source}\\
+\text{source-linked shifted child if required}\\[1mm]
\downarrow\\
\text{Gate 1B closed}\\[2mm]
+\text{source-exhaustive Ford--Maynard packet compiler}\\
\downarrow\\
\text{downstream sieve rerun.}
\end{array}
\]

\subsection{Literature positioning}

The determinant-character bounds of Fouvry--Shparlinski \cite{FouvryShparlinski2022} and the
monomial/reciprocal-monomial character-sum technology of Shparlinski
\cite{Shparlinski2013} are natural candidates only after a literal shifted-source dictionary is
recovered.  Bourgain--Garaev \cite{BourgainGaraev2012} provide reciprocal-set and multilinear
Kloosterman technology that may become relevant to an explicit reciprocal
product dictionary.  Blomer--Pascadi \cite{BlomerPascadi2026} provide strong fixed-modulus bilinear
Kloosterman capacity, but their theorem does not by itself supply the
physical source push-forward norm in \eqref{eq:source-spread}.  Thus the
current first blockers are source-energy and source-definition problems,
not the absence of another generic fixed-multiplier theorem.

\subsection{Status before the high-conductor signed-parent correction}

The following ledger records the positive children established before Late update~IV; it is not the controlling high-conductor status.  The ledger is:
\[
\begin{array}{ll}
\text{native first-}TT^\ast\text{ sector law}
 & \text{proved internally},\\
\text{same-}q\text{ smooth modular-hyperbola discrepancy}
 & \text{analytic pass on proper smooth sector},\\
\text{same-}q\text{ cross-character Gram}
 & \text{analytic pass},\\
\text{cross-}q\text{ no-wrap moving-}D\ L^2
 & \text{proved at operator level},\\
\text{same-}q\text{ diagonal source energy}
 & \text{open},\\
\text{cross-}q\text{ physical source spread}
 & \text{open},\\
\text{actual switched }c_9\text{ dictionary}
 & \text{source open},\\
\text{shifted literal }TT^\ast\text{ source}
 & \text{source open},\\
\text{Gate 1B}
 & \textbf{open},\\
\text{Full Type II / twin-prime infinitude}
 & \textbf{not proved}.
\end{array}
\]




\section{Late update IV: high-conductor complement sign recovery and common-conductor covariance}
\label{app:hicond-common-conductor}

This section records the first high-conductor signed-parent correction.
It is retained for audit history, but its final claim that the common-conductor
covariance is the controlling atomic open is superseded by Late update~V.
It records only identities and reductions supported by the current source
derivation.  In particular, no application of Pascadi's Proposition~3.8 or of
Blomer--Pascadi is claimed here.  The principal architectural correction is
that a repair which genuinely uses the M\"obius sign of the modulus coefficient
cannot close the same-$q$ high-conductor sector in isolation.

\subsection{Banked low/mid-conductor material}

We retain the previously audited same-$q$ collision dictionary and the
low/mid-conductor same-$q$ estimate.  The latter still carries the source-audit
requirement that the primitive large-sieve transcription use the correct
$1/\phi(c)$ weighting and the actual common coefficient sequence.  These
children are not reopened below.

The native proper QK source remains
\[
 \sum_{C_5,x_2,x_3,q}
 \Gamma_5(C_5)f_2(x_2)f_3(x_3)\beta(q)\frac{Y^2}{q^2}
 \sum_{h,k}\Psi_{C_5,x_2,x_3,q}(h,k)
 S(hk\overline B,\pm2;q),
 \qquad
 B=C_5x_2x_3\asymp Y^7,
\]
with $|h|,|k|\ll q/Y$.  The high-conductor correction below undoes only the
outermost Cauchy/character-$L^2$ step needed to recover the signed parent; it
does not replace the physical source by a new model.

\subsection{Primitive conductor complement and exact sign factorisation}

Write
\[
   q=cr,
   \qquad
   c\asymp C,
   \qquad
   C>C_*:=(QY^{11})^{1/3}.
\]
Then
\[
   r\ll \frac QC
   \le \frac Q{C_*}
   =\left(\frac{Q^2}{Y^{11}}\right)^{1/3}
   \ll Y^{5/3+o(1)},
\]
since $Q<Y^8$.  The physical prime slot satisfies
$p>V=Y^{5/2-o(1)}$, hence $p\nmid r$ in the residual band and the physical
prime belongs to the primitive conductor $c$.  Writing $c=pc_0$ and
$d=c_0r$, the literal high-conductor source factorisation is
\begin{equation}
   \boxed{\beta(cr)=\mu(r)\,\varrho_{C,r}(c),}
   \label{eq:beta-complement}
\end{equation}
where
\[
 \varrho_{C,r}(c)
 =
 \sum_{\substack{pc_0=c\\
                  p\text{ and }c_0r\text{ satisfy the physical boxes}}}
   \mu(c_0)\log p.
\]
On the clean squarefree cell, $\mu(c_0)=-\mu(c)$ and therefore
\begin{equation}
  \boxed{
  \beta(cr)=-\mu(c)\mu(r)L_{D,P}(cr),
  }
  \label{eq:beta-clean}
\end{equation}
with
\[
 L_{D,P}(cr)
 =
 \sum_{\substack{p\mid c,\ p\in\mathcal P\\
                  (c/p)r\in\mathcal D}}
 \log p\ge0.
\]
The selector $L_{D,P}(cr)$ is genuinely joint in $(c,r)$; no tensorisation
$L_1(c)L_2(r)$ is used or claimed.  We record
\[
  \boxed{\texttt{BETA-HICOND-COMPLEMENT-FACTOR45: internal proof}.}
\]

\subsection{Induced Gauss square and survival of the complement sign}

Let $\chi\pmod{cr}$ be induced by a primitive character
$\chi^*\pmod c$.  On the proper squarefree sector,
\begin{equation}
   \tau_{cr}(\chi)=\mu(r)\chi^*(r)\tau_c(\chi^*).
   \label{eq:induced-gauss}
\end{equation}
The QK character expansion contains the square of the Gauss sum, so
\[
   \tau_{cr}(\chi)^2
   =\chi^*(r)^2\tau_c(\chi^*)^2.
\]
Combining this with \eqref{eq:beta-complement} gives the exact identity
\begin{equation}
  \boxed{
  \beta(cr)\tau_{cr}(\chi)^2
   =
  \mu(r)\varrho_{C,r}(c)\chi^*(r)^2\tau_c(\chi^*)^2.
  }
  \label{eq:gauss-square-sign}
\end{equation}
Thus one factor $\mu(r)$ survives in the QK square-Gauss parent.  This is the
key distinction from the earlier one-Gauss expression, in which the induced
Gauss factor spends the short-cofactor sign.  We record
\[
 \boxed{\texttt{QK-GAUSS-SQUARE-SIGN-SURVIVES45: internal proof}.}
\]

\subsection{Exact signed parent before the sign-killing Cauchy step}

For the proper unit sector define
\[
 \mathcal A_{c,r,\chi^*}(h,k)
 :=
 \sum_{\substack{B\asymp Y^7\\(B,cr)=1}}
   \mathcal G_{cr}(B;h,k)\chi^*(B),
\]
where $\mathcal G_{cr}$ contains the literal factors
$\Gamma_5(C_5)$, $f_2(x_2)$, $f_3(x_3)$,
$\Psi_{C_5,x_2,x_3,cr}(h,k)$ and $B=C_5x_2x_3$.  Substituting
\eqref{eq:beta-complement} and \eqref{eq:gauss-square-sign} into the one-copy
character formula yields
\begin{equation}
\boxed{
\begin{aligned}
 \mathcal S_{\rm hi}^{\rm preC}
 ={}&
 Y^2\!\sum_{\substack{c>C_*\\cr\asymp Q\\(c,r)=1}}
 \frac{\mu(r)\varrho_{C,r}(c)}{c^2r^2\phi(cr)}
 \sum_{\chi^*\bmod c}^{*}
 \tau_c(\chi^*)^2\chi^*(r)^2  \\
 &\qquad\qquad\qquad\times
 \sum_{h,k}\overline{\chi^*}(bhk)
 \mathcal A_{c,r,\chi^*}(h,k),
\end{aligned}}
\label{eq:hicond-prec}
\end{equation}
with $b=\pm2$ according to the fixed source convention.  On the clean
squarefree cell, \eqref{eq:beta-clean} gives the equivalent form with
$-\mu(c)\mu(r)L_{D,P}(cr)$ in the numerator.  Axes, non-unit,
repeated-prime and shared-prime strata remain in their existing routers.
Hence
\[
  \boxed{\texttt{HICOND-SIGNED-PARENT-RECOVERED: internal proof}.}
\]

\subsection{Why the old SHORT-RJ pre-Cauchy geometry is false}

The positive collision $B_2-B_1=jc$ arises only after opening
$\sum_{\chi^*}|\mathcal A_{c,r,\chi^*}|^2$ and applying character
orthogonality.  That same step replaces the linear source coefficient
$\beta(cr)$ by $|\beta(cr)|^2$ and therefore destroys the complement M\"obius
sign.  Conversely, after undoing that Cauchy step, the signed parent
\eqref{eq:hicond-prec} contains one $B$ variable and no shift $j$.
Consequently the numerical identity $R_*J_*=H^{1/3}$ belongs to the
post-Cauchy positive collision geometry and cannot simultaneously be used as a
pre-Cauchy signed entropy resource.  We therefore record the death
certificate
\[
 \boxed{\texttt{SHORT-RJ-PRECAUCHY-SIMULTANEOUS-DICTIONARY45: false}.}
\]

\subsection{High conductor must recombine same-$q$ and common-conductor cross-$q$}

For any family $T_q$,
\[
 \left|\sum_q\beta(q)T_q\right|^2
 =
 \sum_q|\beta(q)|^2|T_q|^2
 +
 \sum_{q_1\ne q_2}
   \beta(q_1)\overline{\beta(q_2)}
   T_{q_1}\overline{T_{q_2}}.
\]
The same-$q$ part only sees $|\beta(q)|^2$ and therefore cannot use the
M\"obius sign.  Any high-conductor argument which genuinely uses the recovered
sign must retain cross-$q$ covariance.  In the common-conductor sector
$q_i=cr_i$, the sign carrier is $\mu(r_1)\mu(r_2)$ and the same-$q$ case
$r_1=r_2$ is merely its positive diagonal.  Thus
\[
 \boxed{\texttt{HICOND-COMMON-CONDUCTOR-RECOMBINATION45: internal proof}.}
\]
This supersedes the earlier sequential plan ``close same-$q$, then cross-$q$''
for the high-conductor floor.

\subsection{Square-complement collision and the no-wrap threshold}

If one now takes a character $L^2$ step while retaining the
$(r_1,r_2)$ coherence, the primitive-character projector gives, on the proper
unit sector,
\begin{equation}
\boxed{
 \sum_{\chi^*\bmod c}^{*}
 \chi^*(r_1^2B_1)\overline{\chi^*(r_2^2B_2)}
 =
 \sum_{d\mid(c,r_1^2B_1-r_2^2B_2)}
 \phi(d)\mu(c/d).
}
\label{eq:primitive-projector-square}
\end{equation}
Hence the top $d=c$ stratum has the collision
\begin{equation}
   \boxed{r_1^2B_1-r_2^2B_2=j_{\rm sig}c,}
   \label{eq:signed-square-collision}
\end{equation}
not $B_1-B_2=jc$.  If $r_i\asymp R=Q/C$ and $B_i\asymp Y^7$, then
\[
  |j_{\rm sig}|
  \ll \frac{R^2Y^7}{C}
  =\frac{Q^2Y^7}{C^3}.
\]
The corresponding no-wrap threshold is
\begin{equation}
   \boxed{C_{\rm nw}:=(Q^2Y^7)^{1/3}.}
   \label{eq:cnw}
\end{equation}
Moreover
\[
   C_{\rm nw}<Q\quad\Longleftrightarrow\quad Q>Y^7,
\]
so the global phase boundary $\omega=7/9$ reappears for a precise reason:
it is the point at which a signed square-complement no-wrap slice exists.
For $C>C_{\rm nw}$, on the fully proper cross-coprime squarefree stratum,
the top collision reduces to $r_1^2B_1=r_2^2B_2$; under
$(B_1B_2,r_1r_2)=1$, unique factorisation forces $r_1=r_2$ and $B_1=B_2$.
Thus the top projector is positive diagonal in the upper no-wrap slice; any
cancellation must come from a proper-divisor projector or a separately routed
shared-prime stratum.

\subsection{Source-stage firewall against a direct MAM45/HFMV identification}

The recovered parent \eqref{eq:hicond-prec} already lies after
2-dimensional Poisson completion, Kloosterman completion and a
primitive-conductor character projection.  Undoing the final Cauchy step does
not invert these earlier transforms or remove the analytic selector
$\operatorname{cond}(\chi)=c>C_*$.  Therefore the present source does not
supply a literal operator identity from the high-conductor QK parent to the
pre-completion MAM45 or HFMV parents.  We record
\[
 \boxed{\texttt{DIRECT HICOND-QK-PARENT} \to
        \texttt{MAM45/HFMV}: \text{source-stage mismatch}.}
\]
A global MAM45/HFMV theorem could of course control the full physical parent
and hence its high-conductor residual after subtraction of already-controlled
pieces; that is a global theorem landing, not a local dictionary.

\subsection{The exact atomic high-conductor covariance}

Put
\[
 \mathcal Z_{c,r,\chi}(h,k)
 :=
 \overline\chi(bhk)\mathcal A_{c,r,\chi}(h,k).
\]
The exact common-conductor covariance is
\begin{equation}
\boxed{
\begin{aligned}
 \mathcal C_{\rm cc}^{\rm hi}
 :={}&
 Y^4\sum_{c>C_*}
 \sum_{\substack{r_1,r_2\\cr_i\asymp Q}}
 \frac{\mu(r_1)\mu(r_2)
       \varrho_{C,r_1}(c)\overline{\varrho_{C,r_2}(c)}}
      {c^4r_1^2r_2^2\phi(cr_1)\phi(cr_2)}  \\
 &\times
 \sum_{\chi^*\bmod c}^{*}|\tau_c(\chi^*)|^4
 \chi^*\!\left(\frac{r_1^2}{r_2^2}\right)
 \sum_{h,k}
 \mathcal Z_{c,r_1,\chi^*}(h,k)
 \overline{\mathcal Z_{c,r_2,\chi^*}(h,k)}.
\end{aligned}}
\label{eq:cc-hi}
\end{equation}
The required high-conductor theorem is
\begin{equation}
  \boxed{
  \mathcal C_{\rm cc}^{\rm hi}
  \ll_A
  \frac{X^2Y^2}{Q^2}(\log X)^{-A}.
  }
  \label{eq:cc-hi-target}
\end{equation}
This single parent contains the positive same-$q$ floor $r_1=r_2$, the
cross-$q$ terms which carry the only available complement M\"obius sign, the
primitive-projector collision \eqref{eq:signed-square-collision}, and the
actual five-defect QK coefficients.  It is therefore the controlling atomic
high-conductor open.

\subsection{Square-complement adapter: candidate only}

Equation \eqref{eq:gauss-square-sign} shows that the complement phase is
quadratic in $r$.  Formally, if a later source-preserving Kloosterman
transcription produces $\overline r^{\,2}$, the positive-integer map
$s=r^2$ is injective and preserves both $\ell^1$ and $\ell^2$ norms of the
outer weights.  In the high-conductor range,
\[
  \frac{R^2}{C}=\frac{Q^2}{C^3}
  < \frac{Q}{Y^{11}}
  \ll Y^{-3+o(1)}.
\]
This makes $s=r^2$ a genuinely short outer variable relative to the primitive
conductor.  It motivates a fresh audit against Pascadi's variable-family
Kloosterman machinery \cite{Pascadi2025}, but no application is claimed here.
Before such an invocation one must prove the literal $\overline r^{\,2}$
Kloosterman slot, place the joint selector $\varrho_{C,r}(c)$ in a legal
coefficient architecture, expose the genuine multiplicative smooth slots,
verify the relevant exceptional-spectrum hypothesis, and recompute every
geometric term after replacing the outer range $R$ by $R^2$.  These are open
source/analytic tasks.

\subsection{Controlling Gate--1B ledger after the correction}

The current high-conductor ledger is
\[
\begin{array}{ll}
\text{same-$q$ collision dictionary} & \text{banked},\\
\text{same-$q$ low/mid conductor} & \text{banked; weighting audit tag retained},\\
q=cr\text{ complement geometry} & \text{proved internally},\\
\beta(cr)=\mu(r)\varrho_{C,r}(c) & \text{proved internally},\\
\tau_{cr}(\chi)=\mu(r)\chi^*(r)\tau_c(\chi^*) & \text{proved internally},\\
\beta(cr)\tau_{cr}(\chi)^2\text{ retains }\mu(r) & \text{proved internally},\\
\text{signed pre-Cauchy parent} & \text{recovered},\\
\text{old SHORT-RJ pre-Cauchy dictionary} & \textbf{false / retracted},\\
\text{correct signed collision} & r_1^2B_1-r_2^2B_2=j_{\rm sig}c,\\
C_{\rm nw}=(Q^2Y^7)^{1/3} & \text{proved threshold},\\
\text{same-$q$ high conductor as standalone sign repair} & \text{invalid architecture},\\
\text{common-conductor recombination} & \text{proved structurally},\\
\text{direct high-QK to MAM45/HFMV} & \text{source-stage mismatch},\\
\mathcal C_{\rm cc}^{\rm hi}\text{ target} & \textbf{first high-conductor analytic open},\\
\text{square-complement Pascadi adapter} & \text{candidate/open},\\
\text{Gate 1B} & \textbf{open}.
\end{array}
\]
Accordingly this section replaces the old SHORT-RJ and standalone same-$q$
high-conductor architectures.  The subsequent square-complement audit in
Late update~V goes further: it proves that the common-conductor block is not
the full parity-sensitive parent and retires the direct Pascadi-3.8 closure
route.

\section{Late update V: square-complement death certificate and the near-primitive first open}
\label{app:hicond-nearprim}

This section is the controlling Gate--1B update for the high-conductor range.
It records the hostile audit of the common-conductor square-complement adapter.
The outcome is a genuine DAG repair rather than a positive Pascadi splice:
the short-$r$ square norm and a primitive-Kloosterman inversion are exact, but
the naive $r^2$ dictionary is false, both legal Proposition~3.8 lanes are
power-nonclosing, and the common-conductor covariance has already squared away
the conductor M\"obius sign.  The unavoidable first analytic child is therefore
the near-primitive full-conductor five-defect signed packet.

\subsection{Exact short-complement square norm}

Put
\[
 \lambda_{c,r}:=\frac{\mu(r)\varrho_{C,r}(c)}{r^2\phi(cr)}.
\]
Using
\[
 \chi^*\!\left(\frac{r_1^2}{r_2^2}\right)
 =\chi^*(r_1)^2\overline{\chi^*(r_2)^2},
\]
the common-conductor covariance \eqref{eq:cc-hi} is exactly
\begin{equation}
\boxed{
\begin{aligned}
 \mathcal C_{\rm cc}^{\rm hi}
 =Y^4\sum_{c>C_*}\frac1{c^4}
 \sum_{\chi^*\bmod c}^{*}|\tau_c(\chi^*)|^4
 \sum_{h,k}
 \left|
   \sum_{\substack{r\\cr\asymp Q}}
      \lambda_{c,r}\chi^*(r)^2
      \mathcal Z_{c,r,\chi^*}(h,k)
 \right|^2.
\end{aligned}}
\label{eq:cc-short-r-square}
\end{equation}
No estimate is used in \eqref{eq:cc-short-r-square}; it is simply a reassembly
of the $r_1,r_2$ covariance into a short-complement square norm.  We record
\[
 \boxed{\texttt{COMMON-CONDUCTOR-SHORT-r-SQUARE-NORM45: internal proof}.}
\]

\subsection{The genuinely short quantity and the physical aggregate variable}

With $R=Q/C$ and $C>C_*=(QY^{11})^{1/3}$,
\begin{equation}
 \frac{R^2}{C}=\frac{Q^2}{C^3}<\frac{Q}{Y^{11}}
 <Y^{-3+o(1)},
 \qquad
 R^2\ll CY^{-3+o(1)}.
 \label{eq:r2-over-c}
\end{equation}
Thus $r^2$ by itself is a very short integer variable relative to the primitive
conductor.  The decisive correction is that the physical Kloosterman inverse
variable is not $r^2$ alone.  Since $B\asymp Y^7$, the aggregate scale is
\begin{equation}
 \frac{R^2Y^7}{C}=\frac{Q^2Y^7}{C^3},
 \label{eq:r2B-scale}
\end{equation}
which is $<1$ only after $C>C_{\rm nw}=(Q^2Y^7)^{1/3}$.  Hence
\eqref{eq:r2-over-c} does not supply a free short-variable theorem throughout
the high-conductor range.

\subsection{Primitive-Kloosterman inversion}

For squarefree $c$ define the primitive kernel
\[
 \mathfrak P_c(x):=
 \sum_{\chi\bmod c}^{*}\tau_c(\chi)^2\overline\chi(x).
\]
Decomposing every character modulo $c$ by its primitive conductor and using
the induced-Gauss identity gives
\[
 \phi(c)S(1,x;c)
 =\sum_{d\mid c}
   \mathfrak P_d\!\left(x\overline{(c/d)^2};d\right).
\]
M\"obius inversion on the divisor lattice therefore yields the exact identity
\begin{equation}
 \boxed{
 \mathfrak P_c(x)
 =\sum_{d\mid c}\mu(c/d)\phi(d)
 S\!\left(1,x\overline{(c/d)^2};d\right).
 }
 \label{eq:primitive-kloos-inversion}
\end{equation}
For $c=p$ prime this becomes
$\mathfrak P_p(x)=(p-1)S(1,x;p)-1$, as expected from ``all characters minus
principal''.  Substituting the physical
$x=bhk\overline{r^2B}$ and writing $c=de$ gives
\begin{equation}
 \boxed{
 \mathfrak P_c\!\left(bhk\overline{r^2B}\right)
 =\sum_{de=c}\mu(e)\phi(d)
 S\!\left(hk\overline{(er)^2B},b;d\right).
 }
 \label{eq:primitive-kloos-physical}
\end{equation}
Thus even the top $d=c$ term contains the inverse of $r^2B$, not the inverse
of $r^2$ alone.  We record
\[
 \boxed{\texttt{PRIMITIVE-KLOOSTERMAN-INVERSION45: internal proof}.}
\]

\subsection{Why the naive Pascadi-3.8 square-complement dictionary is false}

Pascadi's Proposition~3.8 is a theorem for a kernel of the schematic form
\[
 S(m_1m_2\overline{\mathfrak r},\pm n;\mathfrak s\mathfrak c)
\]
with two genuine multiplicative-convolution slots $m_1,m_2$ and an arithmetic
sequence subject to its Assumption~3.2 \cite{PascadiDistribution2025}.  The physical top
kernel in \eqref{eq:primitive-kloos-physical} is
\[
 \boxed{S(hk\overline{r^2B},b;c),\qquad B=C_5x_2x_3.}
\]
The two multiplicative slots are already forced to be $m_1=h$ and $m_2=k$.
Consequently the five defects and two model factors cannot be reassigned to
those slots without changing the source.  There are only two legal broad
arrangements.

First, scaling $B$ into the second Kloosterman argument gives
\[
 S(hk\overline{r^2B},b;c)
 =S(hk\overline{r^2},b\overline B;c),
\]
so one may take $\mathfrak r=r^2$ and encode
$n\equiv b\overline B\pmod c$ in an inverse-residue coefficient sequence.
This is structurally legal, and the general-sequence form of Assumption~3.2
only inserts the positive norm
\[
 A_{r^2,c}^2=\sum_n|a_{n,r^2,c}|^2.
\]
The theorem therefore requires, rather than proves, the physical inverse-residue
spread.  The literal parameter audit in the present source yields at best
\begin{equation}
 |\mathcal S_{\rm hi}(C)|
 \ll X^{o(1)} C Q^{1/2}A_C.
 \label{eq:p38-inverse-output}
\end{equation}
In the residue-injective upper range $C,Q\ge Y^7$, the natural value
$A_C=Y^{7/2+o(1)}$ leaves a deficit at least $Y^5$.  Hence
\[
 \boxed{\texttt{LEGAL-P38-INVERSE-RESIDUE-LANE45: power-nonclosing}.}
\]

Second, one may put the entire product $r^2B$ into the outer inverse variable.
Then
\[
 \mathfrak R\asymp R^2Y^7=\frac{Q^2Y^7}{C^2},
\]
and the resulting geometric factor and source norm are even larger; the
source audit gives an output of size
\[
 Y^{23/2+o(1)}\sqrt{CQ},
\]
which is far above the physical target $Y^9$.  Hence
\begin{center}
\small\texttt{LEGAL-P38-PRODUCT-OUTER-LANE45}: 
\textbf{EVEN-MORE-NONCLOSING}.
\end{center}
The common-conductor Pascadi closure route is therefore retired.  This is a
capacity and source-decomposition death certificate; it is not a claim that
Proposition~3.8 is incorrect.

\subsection{The deeper sign loss inside the common-conductor covariance}

On the clean cell,
\[
 \beta(cr)=-\mu(c)\mu(r)L_{D,P}(cr).
\]
But within a common-conductor covariance,
\[
 \beta(cr_1)\overline{\beta(cr_2)}
 =\mu(c)^2\mu(r_1)\mu(r_2)L_{D,P}(cr_1)L_{D,P}(cr_2),
\]
so the conductor sign $\mu(c)$ has already been squared away.  The block only
retains the complement pair $\mu(r_1)\mu(r_2)$.  In particular the
near-primitive slice $r_1=r_2=1$ contains no complement sign at all.  Thus
\[
 \boxed{\texttt{COMMON-CONDUCTOR-ONLY-SIGN-REPAIR45: incomplete}.}
\]
The common-conductor square norm remains a legitimate secondary $r>1$ child,
but it cannot be treated as the full parity-breaking high-conductor parent.

\subsection{Representation-loop interpretation}

Equation \eqref{eq:primitive-kloos-inversion} also explains why the
square-complement route does not evade the earlier fixed-modulus death
certificate.  If the primitive conductor $c$ is kept fixed, the quadratic
complement phase remains visible but there is no single complete Kloosterman
sum.  If one expands the primitive projector into complete Kloosterman sums,
one must sum over $d\mid c$ and replace $r$ by $er$ with $e=c/d$; summing all
primitive conductors reconstructs the original full-character QK5 source.
The source energy therefore migrates into the inverse-residue norm or the
outer-weight norm rather than disappearing.

\subsection{The unavoidable near-primitive signed child}

Every $r^2$-based complement argument must confront $r=1$.  After a fixed
polylogarithmic nuclear separation, one near-primitive source packet has the
form
\begin{equation}
\boxed{
\begin{aligned}
 \mathcal S_{\rm NP}^{(\nu)}
 ={}&-Y^2\sum_{c\asymp Q}
 \frac{\mu(c)L_{D,P}(c)}{c^2\phi(c)}
 \sum_{\chi^*\bmod c}^{*}\tau_c(\chi^*)^2
 \prod_{i=1}^5D_{i,\nu,c}(\chi^*)
 F_{2,\nu,c}(\chi^*)F_{3,\nu,c}(\chi^*)\\
 &\hspace{38mm}\times
 \overline{H_{\nu,c}(\chi^*)K_{\nu,c}(\chi^*)}.
\end{aligned}}
\label{eq:nearprim-signed}
\end{equation}
Here the conductor M\"obius sign $\mu(c)$ and all five centred defect factors
remain linear; there is no complement variable available for square-complement
averaging.  The first unavoidable high-conductor theorem is therefore
\begin{equation}
 \boxed{
 \mathcal S_{\rm NP}^{(\nu)}\ll_A X(\log X)^{-A}.
 }
 \label{eq:nearprim-target}
\end{equation}
Equivalently, at the pre-completion source stage one must bound the signed
native determinant packet
\begin{equation}
 \boxed{
 \sum_{\substack{B,x_1,x_4,c,\ell\\Bx_1x_4-c\ell=-2}}
 \Gamma_5(C_5)f_2(x_2)f_3(x_3)
 \mu(c)L_{D,P}(c)W(\cdots)
 \ll_A X(\log X)^{-A}.
 }
 \label{eq:nearprim-precompletion}
\end{equation}
For $r=1$, the top $d=c$ term of
\eqref{eq:primitive-kloos-inversion} is precisely
$\phi(c)S(hk\overline B,b;c)$, the original QK5 complete kernel.  Thus the
finite/source inversion back to the native determinant shell
$Bx_1x_4-c\ell=-2$, while preserving
$\beta(c)=-\mu(c)L_{D,P}(c)$, is an exact source adapter.  It is not an
estimate.  We record
\[
 \boxed{\texttt{NEARPRIM-TOP-KLOOSTERMAN-TO-SIGNED-QK5/MAM45-ADAPTER: pass}.}
\]

\subsection{Revised Gate--1B high-conductor DAG}

The controlling high-conductor path is now
\[
\boxed{
\begin{array}{c}
\text{low/mid conductor bank}
+\text{ cross-character bank}
+\text{ cross-$q$ no-wrap operator bank}
\\[1mm]\downarrow\\[-1mm]
\textbf{HIGH CONDUCTOR}
\quad
\begin{cases}
 r=1:&\boxed{\texttt{NEARPRIM-FULL-c-FIVEDEFECT-SIGNED45}},\\
 r>1:&\text{square-character complement family},
\end{cases}
\\[2mm]\downarrow\\
\text{full }(c_1,r_1)\times(c_2,r_2)\text{ source reassembly}
\\\downarrow\\
\text{QK56 parent}
\\\downarrow\\
\text{exhaustiveness / shifted-child test}
\\\downarrow\\
\texttt{GATE1B CLOSED}.
\end{array}}
\]
The common-conductor covariance is retained only as a secondary $r>1$ child.
The present first analytic open is
\texttt{NEARPRIM-FULL-c-FIVEDEFECT-SIGNED45}.

\subsection{Controlling ledger after the square-complement audit}

\[
\begin{array}{ll}
\text{common-conductor short-$r$ square norm} & \text{proved exactly},\\
R^2/C\ll Y^{-3+o(1)} & \text{proved},\\
\text{primitive-Kloosterman inversion} & \text{proved exactly},\\
\text{actual inverse variable} & r^2B\text{, not }r^2,\\
\text{naive }r^2\to\text{Pascadi-3.8 dictionary} & \textbf{false},\\
\text{legal inverse-residue P3.8 lane} & \textbf{power-nonclosing},\\
\text{legal product-outer P3.8 lane} & \textbf{even-more-nonclosing},\\
\text{common-conductor }\mu(c) & \text{squared away},\\
r=1\text{ near-primitive child} & \textbf{unresolved signed floor},\\
\text{common-conductor P3.8 closure} & \textbf{retired},\\
\text{new first high-conductor open} & \boxed{\texttt{NEARPRIM-FULL-c-FIVEDEFECT-SIGNED45}},\\
\text{Gate 1B} & \textbf{open}.
\end{array}
\]


\section{Late update VI: source-minimal native \texorpdfstring{$4$M$\mid5$D}{4M|5D} MAM45 reduction}
\label{app:mam45-native}

\begin{status}[Controlling Gate--1B correction]
This section supersedes only the \emph{first-open designation} of Late
update~V.  The primitive-projector inversion, the square-complement death
certificate, and the common-conductor sign-loss analysis remain banked.  The
new correction is that the near-primitive primitive-character packet is not
the source-minimal theorem: its top divisor is an exact child of the full
native signed MAM45 source, while its proper-divisor terms are triangular
$r>1$ children.  Gate~1B remains open.
\end{status}

\subsection{Literal order-five source and the global sign}
For a legitimate smooth/nuclear packet $\nu$, the raw order-five source has the
form
\begin{equation}
\begin{aligned}
 \mathcal S_{\rm raw}^{(\nu)}={}&
 \sum_{\substack{r_1,\ldots,r_5,\,x_1,\ldots,x_4,\,c,\ell\\
 (r_1\cdots r_5)(x_1\cdots x_4)+2=c\ell}}
 \left(\prod_{i=1}^{5}\delta_i(r_i)\right)
 \left(\prod_{j=1}^{4}f_j(x_j)\right)
 \beta_{D,P}(c)W_R(\ell)\,
 \Omega_\nu(\mathbf r,\mathbf x,c,\ell).
\end{aligned}
\label{eq:d14-raw-source}
\end{equation}
The fresh physical centering is
\[
 f_\lambda(m)=\frac{W_\lambda(m)}{\log m},\qquad
 \delta_\lambda(m)=(\Lambda(m)-1)\frac{W_\lambda(m)}{\log m},
\]
with the prime-power discrepancy sent to the already banked fixed-power
router.  On the clean squarefree switched cell,
\begin{equation}
 \beta_{D,P}(c)
 =\sum_{\substack{dp=c\\d\in\mathcal D,\ p\in\mathcal P}}
   \mu(d)\log p
 =-\mu(c)L_{D,P}(c).
\label{eq:d14-beta-sign}
\end{equation}
Thus any earlier schematic convention with $+\mu(c)L_{D,P}(c)$ differs from
the literal raw source by a global minus sign.  This is immaterial for an
absolute bound but load-bearing for the source dictionary.

\subsection{Exact scalarisation into the native \texorpdfstring{$4$M$\mid5$D}{4M|5D} parent}
Set
\begin{equation}
 u=x_1x_2x_3x_4\asymp Y^4=:U,
 \qquad
 v=r_1r_2r_3r_4r_5\asymp Y^5=:V.
\label{eq:d14-uv}
\end{equation}
Then the native determinant shell is exactly
\begin{equation}
 \boxed{uv+2=c\ell.}
\label{eq:d14-native-shell}
\end{equation}
For a separated labelled packet write
\[
 a_\nu(u)
 =\sum_{x_1x_2x_3x_4=u}
   \prod_{j=1}^{4}f_j(x_j)\,\Omega_{\nu,U}(\mathbf x),
\]
and
\[
 b_\nu(v)
 =\sum_{r_1\cdots r_5=v}
   \prod_{i=1}^{5}\delta_i(r_i)\,\Omega_{\nu,V}(\mathbf r).
\]
Labelled dyadic boxes contribute only fixed or $X^{o(1)}$ factorisation
multiplicity.  The top native packet becomes
\begin{equation}
 \mathcal S_{\rm top}^{(\nu)}
 =\sum_{\substack{u\sim U,\ v\sim V,\ c\sim Q,\ \ell\sim R\\
 uv+2=c\ell}}
 a_\nu(u)b_\nu(v)\beta_{D,P}(c)W_R(\ell).
\label{eq:d14-mam-c}
\end{equation}
Opening $c=dp$ gives the exact signed parent
\begin{equation}
 \boxed{
 \mathcal S_{\rm top}^{(\nu)}
 =\sum_{\substack{u,v,d,p,\ell\\uv+2=dp\ell}}
 a_\nu(u)b_\nu(v)\mu(d)\log p\,W_R(\ell).
 }
\label{eq:d14-mam-dp}
\end{equation}
This is precisely the source-minimal scalar MAM45 architecture: four model
coordinates are absorbed in $a_\nu$, five centred defect coordinates in
$b_\nu$, and the M\"obius sign remains \emph{linear} in the pre-Cauchy
$d$-sum.  We therefore bank
\[
 \boxed{\texttt{NEARPRIM-TOP-TO-MAM45-4M5D45: PASS}.}
\]
The analytic target itself remains
\begin{equation}
 \boxed{
 \mathcal S_{\rm MAM}^{(\nu)}
 \ll_A X(\log X)^{-A}.
 }
\label{eq:d14-mam-target}
\end{equation}
We denote this theorem target by
\[
 \boxed{\texttt{MAM45-4M5D-SCALAR}.}
\]

\subsection{Why the full primitive packet is not the native top term}
The primitive projector obeys the previously banked inversion
\[
 \mathfrak P_c(x)
 =\sum_{d\mid c}\mu(c/d)\phi(d)
 S\!\left(1,x\overline{(c/d)^2};d\right).
\]
Only the top divisor $d=c$ inverts to the native packet
\eqref{eq:d14-mam-c}.  Terms $d<c$ introduce the effective complement
$e=c/d>1$, i.e. the triangular $r>1$ conductor children.  Hence
\[
 \boxed{
 \text{primitive near-$c$ packet}
 =\text{native MAM top}
 +\text{triangular $r>1$ divisor children}.
 }
\]
Accordingly
\[
 \boxed{\texttt{FULL-PRIMITIVE-PACKET-(52)=(54): NOT LITERAL}.}
\]
This is a useful weakening: proving the global native MAM45 parent preserves
cancellation among the children and is not required to make every primitive
component small separately.

\subsection{Centred-source normalisation remains a finite/source pin}
Equation~\eqref{eq:d14-mam-dp} is the exact nonzero-determinant raw source.
For the fully centred identity one must still match, literally and in one
normalisation convention,
\[
 E(c),\qquad Z_E(c),\qquad \kappa_4,\qquad
 \text{the Poisson zero mode},\qquad \text{the main-term convention}.
\]
The source audit therefore records
\[
 \boxed{\texttt{NEARPRIM-RAW-MAM45-DICTIONARY: PASS}},
\]
and
\[
 \boxed{\texttt{NEARPRIM-CENTERED-MAM45-DICTIONARY: SOURCE PIN}},
\]
where the remaining pin is exactly
$E/Z_E/\kappa_4$/Poisson-zero-mode normalisation.
This obligation is source/reassembly bookkeeping; it is not a new Kloosterman
estimate.

\subsection{HFMV is a deterministic sufficient theorem, not a banked analytic input}
Define the literal weighted discrepancy
\[
 \mathcal B_\nu(u)
 :=\sum_{d\sim D}\sum_{p\sim P}\mu(d)\log p
   \sum_{v\sim V}b_\nu(v)\,\rho^W_{dp}(uv+2),
\]
where $\rho^W_{dp}$ includes $dp\mid uv+2$, the induced
$\ell=(uv+2)/(dp)\sim R$, the weight $W_R(\ell)$, and the same centred main
term as the scalar source.  Then
\[
 \mathcal S_{\rm MAM}^{(\nu)}=\sum_{u\sim U}a_\nu(u)\mathcal B_\nu(u).
\]
The labelled four-model convolution gives
\[
 \sum_{u\sim U}|a_\nu(u)|^2\ll U(\log X)^{O(1)}.
\]
Consequently the square-function hypothesis
\begin{equation}
 \sum_{u\sim U}|\mathcal B_\nu(u)|^2
 \ll_B UV^2(\log X)^{-B}
\label{eq:d14-hfmv}
\end{equation}
implies \eqref{eq:d14-mam-target} by one Cauchy inequality, with
$B\ge 2A+O(1)$.  Thus
\[
 \boxed{\texttt{HFMV45}\Longrightarrow\texttt{MAM45-4M5D-SCALAR}}
\]
is a deterministic pass.  The analytic hypothesis \eqref{eq:d14-hfmv}
remains open: one-principal pieces are routed, while the mixed
both-nonprincipal core is not closed.

\subsection{A direct one-defect Siegel--Walfisz escape is impossible}
Fix the other eight source coordinates.  With
\[
 A=r_2r_3r_4r_5x_1x_2x_3x_4\asymp Y^8,
\]
the native shell in $(r_1,\ell)$ is $Ar_1+2=c\ell$.  On the unit sector,
\[
 \ell\equiv2\overline c\pmod A.
\]
But $\ell\sim R=X/Q\le Y^{5/2+o(1)}\ll A$, so a dyadic $\ell$ interval
contains at most one point in this residue class.  There is therefore no long
single-defect average after freezing the other eight coordinates:
\[
 \boxed{\texttt{NP-ONEDEFECT-DIRECT-SW45: NO AVERAGING / FAIL}.}
\]
Any use of defect centering must retain a genuinely joint family.

\subsection{Exact determinant-line spectral reduction}
Fix a clean pair $(u,\ell)$ and a base solution $uv_0+2=\ell z_0$.
All solutions on the determinant line are
\[
 v_t=v_0+\ell t,\qquad z_t=z_0+ut,\qquad |t|\ll H,
 \qquad H=\frac QU=\frac VR.
\]
For $M\asymp H$ define
\[
 \mathcal B_{u,\ell}(k)
 =\sum_tW_1(t/H)b_\nu(v_0+\ell t)e_M(-kt),
\]
\[
 \mathcal Z_{u,\ell}(k)
 =\sum_tW_2(t/H)\beta(z_0+ut)e_M(-kt),
\]
with $W_1\overline{W_2}=W$.  Finite Fourier inversion gives exactly
\begin{equation}
 \boxed{
 \mathcal M_{u,\ell}
 =\frac1M\sum_{k\bmod M}
 \mathcal B_{u,\ell}(k)\overline{\mathcal Z_{u,\ell}(k)}.
 }
\label{eq:d14-line-fourier}
\end{equation}
We bank
\[
 \boxed{\texttt{NP-LINE-SPECTRAL-REDUCTION45: PASS}.}
\]
A convenient stronger sufficient target is the joint overlap
\begin{equation}
 \boxed{
 \sum_{u\sim U}\sum_{\ell\sim R}\sum_{k\bmod M}
 |\mathcal B^{(5\delta)}_{u,\ell}(k)|^2
 |\mathcal Z^{(\mu*\Lambda)}_{u,\ell}(k)|^2
 \ll_A MQV(\log X)^{-A}.
 }
\label{eq:d14-line-overlap}
\end{equation}
This target is stronger than scalar MAM45 and is therefore recorded only as a
sufficient route, not an equivalent reformulation.

\subsection{Why the presently banked separate uniformities are power-nonclosing}
The coefficient $\beta=\mu_D*\Lambda_P$ already has arbitrary-log additive
Fourier uniformity and the corresponding additive $U^2$ energy.  In the
post-Cauchy determinant projector, however, the $\beta$ shift is $uh$ while
the five-defect shift is $\ell h$: the same $h$ synchronises both sides.  A
Cauchy split which treats the two energies independently pays the family-size
tax
\[
 \sqrt{UR}=X^{13/18-\omega/2}.
\]
Across the hard band this ranges from $X^{13/36}$ to $X^{5/18}$.  Even if one
credits the full banked five-prime reciprocal saving in both defect copies,
the separated route still leaves a fixed-power deficit ranging from
$X^{1/3-o(1)}$ at the lower endpoint to $X^{1/4-o(1)}$ near the upper endpoint.
Therefore
\[
 \boxed{\texttt{CURRENT-BETA-U2+FIVEDEFECT-BANKS: FIXED-POWER NONCLOSING}.}
\]
The missing resource is joint nonalignment on the \emph{same determinant-line
frequency}; the controlling gap is not merely a loss of two logarithms.

\subsection{Post-Cauchy projector and remaining viable routes}
After one Cauchy step the exact off-diagonal determinant relation is
\[
 z_1v_2-z_2v_1=2h.
\]
The associated five-defect projector gives the scalar quadratic child
\[
 \boxed{\texttt{NP-BETA-BPROJECTOR-GLOBAL45}},
\]
which remains open.  Together with \eqref{eq:d14-line-overlap} and HFMV, this
gives three honest analytic lanes for the source-minimal parent:
\begin{enumerate}[label=(\roman*)]
 \item direct joint line-frequency cancellation for
       \eqref{eq:d14-line-overlap};
 \item the mixed both-nonprincipal core of HFMV;
 \item a source-specific expansion or inverse-rigidity argument for the
       global beta-Bprojector child defined above.
\end{enumerate}
A full-face $1*(\mu*\Lambda)=\Lambda$ recombination is algebraically valid but
does not presently provide a fourth analytic escape: source-face completeness
is not proved, centred transfer leaves an $E(c)$ residual, and the hard-band
$q=n,\ell=1$ face contains a parity-hard shifted M\"obius correlation.

\subsection{Controlling Gate--1B ledger and DAG}
The latest strict ledger is
\begin{center}
\small
\begin{tabular}{p{0.60\textwidth}p{0.28\textwidth}}
\toprule
Item & Status\\
\midrule
\texttt{NEARPRIM-TOP-TO-MAM45-4M5D45} & \textbf{PASS}\\
\texttt{FULL-PRIMITIVE-PACKET-(52)=(54)} & \textbf{NOT LITERAL}\\
\texttt{NEARPRIM-RAW-MAM45-DICTIONARY} & \textbf{PASS}\\
\texttt{NEARPRIM-CENTERED-MAM45-DICTIONARY} & \textbf{SOURCE PIN}\\
\texttt{HFMV45} $\Rightarrow$ \texttt{MAM45-4M5D-SCALAR} & \textbf{DETERMINISTIC PASS}\\
\texttt{HFMV45 analytic input} & \textbf{OPEN}\\
\texttt{NP-ONEDEFECT-DIRECT-SW45} & \textbf{FAIL}\\
\texttt{NP-LINE-SPECTRAL-REDUCTION45} & \textbf{PASS}\\
\texttt{NP-LINE-SPECTRAL-OVERLAP45} & \textbf{OPEN}\\
\texttt{NP-BETA-BPROJECTOR-GLOBAL45} & \textbf{OPEN}\\
\texttt{CURRENT-BETA-U2+FIVEDEFECT-BANKS} & \textbf{POWER-NONCLOSING}\\
\texttt{MAM45-4M5D-SCALAR} & \textbf{FIRST SOURCE-MINIMAL ANALYTIC OPEN}\\
\texttt{GATE1B} & \textbf{OPEN}\\
\bottomrule
\end{tabular}
\end{center}
The shortest honest dependency graph is
\[
\boxed{
\begin{array}{c}
\text{full native }4\text{M}\mid5\text{D MAM45 scalar}\\
\downarrow\\
\begin{cases}
\text{joint line-frequency cancellation},\\
\text{or HFMV mixed-core theorem},\\
\text{or NP-BETA-BPROJECTOR source-specific rigidity}
\end{cases}\\
\downarrow\\
\text{all }r=1\text{ and }r>1\text{ conductor children collectively controlled}\\
\downarrow\\
\text{QK56 parent + exhaustiveness + finite reassembly}\\
\downarrow\\
\texttt{GATE1B CLOSED}.
\end{array}}
\]
This is a research dependency graph, not a proof of its terminal node.  Full
Type-II reassembly, the downstream Ford--Maynard source adapter, and twin-prime
infinitude remain open/conditional.


\section{Late update VII: pure-five-defect coherent projector and entropy-gap router}
\label{app:pure5-dp}

\begin{status}[Historical Gate--1B update; superseded by Late update VIII]
This section superseded the broad statement that
\texttt{MAM45-4M5D-SCALAR} itself was the first unresolved object in the
stipulated source order.  Late update~VIII now supersedes this section as the
controlling status.  The native MAM45 parent remains the correct global
architecture, but the clean $|J|=5$ face has now been reduced to a smaller
source-exact signed projector.  Nothing below proves that projector, and no
lower-defect face is inferred before the $|J|=5$ stopping rule is discharged.
\end{status}

\subsection{Literal pure-five-defect packet}
Write
\[
 Q=Y^\vartheta,\qquad \frac{13}{2}\le\vartheta<8,\qquad X=Y^9,
\]
\[
 U=Y^4,\qquad V=Y^5,\qquad R=X/Q=Y^{9-\vartheta},
\]
and decompose the switched modulus as
\[
 d\asymp D,\qquad p\asymp P,\qquad DP\asymp Q,
 \qquad P>P_0=Y^{5/2-o(1)}.
\]
For the pure five-defect assignment define
\[
 a_4(u)=\sum_{x_1x_2x_3x_4=u}\prod_{j=1}^4 f_j(x_j),
 \qquad
 b_5(v)=\sum_{r_1\cdots r_5=v}\prod_{i=1}^5\delta_i(r_i).
\]
The clean nonzero source is
\begin{equation}
 \boxed{
 \mathcal S_5(D,P)=
 \sum_{\substack{u\sim U,\ v\sim V,\ d\sim D,\ p\sim P,\ \ell\sim R\\
                   uv+2=dp\ell}}
 a_4(u)b_5(v)\mu(d)\log p\,W_R(\ell).
 }
 \label{eq:d15-pure5-source}
\end{equation}
This is a source statement, not an estimate.  On the clean odd sector,
$uv+2=dp\ell$ forces
\[
 (u,dp)=(v,dp)=(u,\ell)=(v,\ell)=1.
\]
Furthermore a squarefree $q=dp<Y^8$ has only $O(1)$ admissible distinguished
prime divisors $p>P_0$: indeed
\[
 \#\{p\mid q:p>P_0\}
 \le \frac{\log q}{\log P_0}<\frac{8}{5/2}+o(1)<4.
\]
We therefore bank
\[
 \boxed{\texttt{PURE5-RAW-SOURCE45: PASS}},\qquad
 \boxed{\texttt{PURE5-DP-FACTORIZATION-MULTIPLICITY45: O(1)}},
\]
and the clean odd coprimality certificate.  The literal comparison/main-term
normalisation is not present in the supplied source and remains
\[
 \boxed{\texttt{PURE5-COMPARISON-MAINTERM-PIN: SOURCE OPEN}}.
\]

\subsection{Determinant-line entropy trichotomy}
Set
\[
 H:=\frac QU=\frac VR=Y^{\vartheta-4},
 \qquad Y^{5/2}\le H<U=Y^4.
\]
Fixing $(d,u,\ell)$, all solutions of $d\ell p-uv=2$ lie on
\[
 p=p_0+ut,\qquad v=v_0+d\ell t,
\]
so the prime-line length is
\[
 L_p\asymp \frac PU.
\]
Fixing $(p,u,\ell)$ instead gives
\[
 d=d_0+ut,\qquad v=v_0+p\ell t,
\]
with M\"obius-line length
\[
 L_d\asymp \frac DU=\frac HP.
\]
Hence
\begin{equation}
 \boxed{L_pL_d=\frac PU\frac HP=\frac HU\le1.}
 \label{eq:d15-entropy-product}
\end{equation}
Equivalently,
\[
\begin{array}{c|c|c}
P\text{-range}&L_p=P/U&L_d=H/P\\ \hline
P<H&<1&>1\\
H\le P\le U&\le1&\le1\\
P>U&>1&<1.
\end{array}
\]
Thus
\[
 \boxed{\texttt{DP-LINE-ENTROPY-TRICHOTOMY45: PASS}}.
\]
The central interval $H\le P\le U$ is an entropy gap for one-dimensional
averaging: neither the prime lane nor the M\"obius lane is long.  In the two
outer bands the long lane is still weighted by the affine pullback of $b_5$,
so ordinary prime-distribution or bare Davenport estimates do not apply
without solving the same five-defect correlation.  We retain the death
certificates
\[
 \boxed{\texttt{PURE5-PRIME-P1-DISPERSION45: SOURCE-COEFFICIENT FAIL}},
\]
\[
 \boxed{\texttt{PURE5-MOBIUS-d-FREQUENCY45: REPRESENTATION LOOP / NONCLOSING}}.
\]

\subsection{Joint line spectrum and inverse certificate}
For fixed clean $(u,\ell)$ choose $uv_0+2=\ell z_0$.  Every point on the
line is
\[
 v_t=v_0+\ell t,\qquad z_t=z_0+ut,\qquad |t|\ll H.
\]
With
\[
 \beta_{D,P}(z)=\sum_{\substack{dp=z\\d\sim D,\ p\sim P}}\mu(d)\log p,
\]
define the common-frequency transforms
\[
 \mathcal B_{u,\ell}(k)
 =\sum_tb_5(v_0+\ell t)W_1(t/H)e_H(-kt),
\]
\[
 \mathcal Z^{D,P}_{u,\ell}(k)
 =\sum_t\beta_{D,P}(z_0+ut)W_2(t/H)e_H(-kt).
\]
Finite Fourier inversion gives
\begin{equation}
 \boxed{
 \mathcal M_{u,\ell}=\frac1H\sum_{k\bmod H}
 \mathcal B_{u,\ell}(k)\overline{\mathcal Z^{D,P}_{u,\ell}(k)}.
 }
 \label{eq:d15-line-dft}
\end{equation}
Opening both transforms displays the simultaneous phases
\[
 \frac{r_1\cdots r_5-v_0}{\ell},\qquad
 \frac{dp-z_0}{u}
\]
at the same frequency $k/H$.  Hence
\[
 \boxed{\texttt{PURE5-JOINT-LINE-PHASE-DICTIONARY45: PASS}}.
\]
Dyadic pigeonholing of simultaneous large values also yields the exact inverse
saturation certificate recorded in the source audit; it does not by itself
supply an upper bound.  We bank
\[
 \boxed{\texttt{PURE5-JOINT-LARGEVALUES-SAT45: PASS}}.
\]

\subsection{The coherent B-projector}
Apply Cauchy only on the model $u$ side.  The diagonal has fixed-power room;
the source-exact off-diagonal is
\begin{equation}
\boxed{
\begin{aligned}
 \mathcal O_{5;D,P}:=
 \sum_{\substack{u\sim U,\ \ell\sim R\\
                 d_1,d_2\sim D,\ p_1,p_2\sim P\\h\ne0}}
 &\mu(d_1)\mu(d_2)(\log p_1)(\log p_2)\\
 &\times b_5(v_1)\overline{b_5(v_2)}\,\mathcal W_\nu,
\end{aligned}}
\label{eq:d15-bprojector}
\end{equation}
subject to
\begin{equation}
 \boxed{d_2p_2-d_1p_1=uh,}
 \qquad
 \boxed{v_i=\frac{\ell d_ip_i-2}{u}\asymp V,}
 \label{eq:d15-coupled-shifts}
\end{equation}
and therefore
\[
 \boxed{v_2-v_1=\ell h},\qquad
 \boxed{d_1p_1v_2-d_2p_2v_1=2h}.
\]
The source-minimal pure-five-defect target is
\begin{equation}
 \boxed{|\mathcal O_{5;D,P}|\ll_A QV(\log X)^{-A}}
 \label{eq:d15-bprojector-target}
\end{equation}
uniformly over $DP\asymp Q$ and $P>P_0$.  We denote
\eqref{eq:d15-bprojector-target} by
\[
 \boxed{\texttt{PURE5-DP-COHERENT-BPROJECTOR45}}.
\]
Its status is \textbf{OPEN}.  It is weaker than the HFMV square-function
because it preserves the global signed scalar sum instead of taking an
absolute value for each $(u,h)$.

\subsection{Equal-factor router and the fully-crossed central band}
If $p_1=p_2=p$, then $p(d_2-d_1)=uh$.  Clean coprimality gives $p\mid h$, so
for $P>H$ the off-diagonal common-$p$ stratum is empty; when $P\le H$ it has
the exact one-parameter form
\[
 h=pj,\qquad d_2-d_1=uj,\qquad |j|\ll H/P.
\]
Similarly, if $d_1=d_2=d$, then $d(p_2-p_1)=uh$ and $d\mid h$.  Since
$D>H$ is equivalent to $P<U$, common-$d$ off-diagonal points are empty for
$P<U$; when $P\ge U$ they satisfy
\[
 h=dj,\qquad p_2-p_1=uj,\qquad |j|\ll P/U.
\]
Thus
\[
\begin{array}{c|c}
P\text{-range}&\text{possible off-diagonal strata}\\ \hline
P<H&\text{common-$p$ + fully crossed}\\
H<P<U&\text{fully crossed only}\\
P>U&\text{common-$d$ + fully crossed}.
\end{array}
\]
We bank
\[
 \boxed{\texttt{EQUAL-LARGE-FACTOR-OFFDIAG-ROUTER45: PASS}}.
\]
In particular, the central band obeys
\[
 \boxed{H<P<U\Longrightarrow d_1\ne d_2,\quad p_1\ne p_2.}
\]

\subsection{New elementary prime-divisor rigidity}
The fully-crossed equation can be rewritten exactly as
\begin{equation}
 \boxed{d_2p_2=d_1p_1+uh.}
 \label{eq:d15-divisor-rigidity}
\end{equation}
For fixed $(d_1,p_1,u,h)$ put $N=d_1p_1+uh$.  On the source ranges
$d_1p_1\asymp Q$ and $|uh|\ll UH=Q$, hence $0<N\ll Q$.  Every admissible
$p_2\asymp P$ is a prime divisor of $N$.  Since $P>P_0=Y^{5/2-o(1)}$ and
$Q<Y^8$,
\[
 \#\{p_2\asymp P:p_2\mid N\}
 \le \frac{\log N}{\log P}
 \le \frac{8+o(1)}{5/2-o(1)}=O(1).
\]
Once $p_2$ is selected,
\[
 \boxed{d_2=\frac{d_1p_1+uh}{p_2}}
\]
is determined and is merely filtered by the $D$-box and clean-source
conditions.  Equivalently, for fixed $(d_2,u,h,p_1)$,
\[
 p_2\equiv uh\,\overline{d_2}\pmod{p_1},
\]
and a dyadic $p_2$ interval of size $\asymp P\asymp p_1$ contains only $O(1)$
integers in this residue class.  We therefore bank the finite certificate
\[
 \boxed{\texttt{FULLYCROSSED-PRIME-DIVISOR-RIGIDITY45: PASS}}.
\]
This collapses the prime-pair entropy but does \emph{not} prove cancellation
in \eqref{eq:d15-bprojector-target}: the remaining $d_1,p_1,u,h,\ell$ family
still carries both M\"obius signs and the coupled five-defect coefficient.

\subsection{Next analytic attack and the two wings}
The cleanest unresolved band is therefore the central fully-crossed child
\[
 \boxed{\texttt{PURE5-DP-FULLYCROSSED-ENTROPYGAP45: NEXT / OPEN}},
\]
meaning \eqref{eq:d15-bprojector-target} restricted to
\[
 H<P<U,\qquad d_1\ne d_2,\qquad p_1\ne p_2.
\]
This is a deliberately narrower attack target, not a claim that the two wings
are already harmless.  After the central child one must still treat
\[
 P<H:\ \text{common-$p$ wing},\qquad
 P>U:\ \text{common-$d$ wing},
\]
in addition to the fully-crossed pieces there.  Their exact one-parameter
routers reduce dimension, but coefficient-blind counting is nonclosing.

The current published Kloosterman literature does not by itself instantiate
this source-specific projector.  Blomer--Pascadi prove strong all-moduli
bilinear Kloosterman bounds, including a $c^{-1/32}$ saving in the critical
square-root range \cite{BlomerPascadi2026}; Pascadi's distribution paper
provides flexible Kloosterman/large-sieve architecture \cite{PascadiDistribution2025}.
These are backend-capacity references, not automatic providers for the
coupled signed projector above.

\subsection{Major arcs, Poisson tail, and stopping-rule ledger}
For
\[
 \widehat\beta_{D,P}(\alpha)
 =\sum_{p\sim P}\log p\sum_{d\sim D}\mu(d)e(\alpha pd),
\]
Davenport-type M\"obius Fourier cancellation closes the large-$D$ major-arc
piece once $D$ exceeds a suitable logarithmic threshold.  It does not close
$D=O(1)$: at $D=1$ the zero frequency is of size $\asymp P$.  Thus the safe
status is
\[
 \boxed{\texttt{PURE5-MAJORARC-LARGED45: PASS}}.
\]
\[
 \boxed{\texttt{PURE5-MAJORARC-SMALLD45: MAINTERM/SOURCE PIN OPEN}}.
\]
Likewise the inverse-Poisson truncation tail is analytically routine under the
stated rapid-decay derivative bounds, but the derivative provenance has not
been printed in the supplied source:
\[
 \boxed{\texttt{D2-RAPID-DECAY-PIN: SOURCE OPEN / ANALYTICALLY ROUTINE}}.
\]

The controlling ledger is therefore
\begin{center}
\small
\begin{tabular}{p{0.61\textwidth}p{0.27\textwidth}}
\toprule
Item & Status\\
\midrule
\texttt{PURE5-RAW-SOURCE45} & \textbf{PASS}\\
\texttt{DP-LINE-ENTROPY-TRICHOTOMY45} & \textbf{PASS}\\
\texttt{PURE5-JOINT-LINE-PHASE-DICTIONARY45} & \textbf{PASS}\\
\texttt{PURE5-JOINT-LARGEVALUES-SAT45} & \textbf{PASS}\\
\texttt{EQUAL-LARGE-FACTOR-OFFDIAG-ROUTER45} & \textbf{PASS}\\
\texttt{FULLYCROSSED-PRIME-DIVISOR-RIGIDITY45} & \textbf{PASS}\\
\texttt{PURE5-DP-FULLYCROSSED-ENTROPYGAP45} & \textbf{NEXT / OPEN}\\
\texttt{PURE5-DP-COHERENT-BPROJECTOR45} & \textbf{OPEN}\\
\texttt{PURE5-COMPARISON-MAINTERM-PIN} & \textbf{SOURCE OPEN}\\
\texttt{HFMV-BOTH-NONPRINCIPAL-CORE45} & \textbf{OPEN}\\
\texttt{PURE5-DP-SIGNED45} & \textbf{REDUCED -- NOT CLOSED}\\
$|J|=4,3,2,1$ defect faces & \textbf{NOT YET INFERRED}\\
$r>1$ square-character family & \textbf{OPEN SECONDARY}\\
CSTAR--CNW transition strip & \textbf{OPEN}\\
QK56 full parent / exhaustiveness & \textbf{OPEN}\\
\texttt{GATE1B} & \textbf{OPEN}\\
\bottomrule
\end{tabular}
\end{center}
At the time of this audit, the first exact residual under the stipulated
source stopping order was
\[
 \boxed{\texttt{PURE5-DP-COHERENT-BPROJECTOR45}.}
\]
Late update~VIII supersedes this headline by reducing the fully-crossed
central child to a double-GCD B\'ezout line and then to
\texttt{RANKONE-BEZOUT-AFFINE-ELLIOTT5D45}.  The present section is retained
as audit history.


\section{Late update VIII: double-GCD B\'ezout line and rank-one inverse frontier}
\label{app:rankone-bezout}

\begin{status}[Controlling Gate--1B update]
This section supersedes the Draft~15 claim that
\texttt{PURE5-DP-COHERENT-BPROJECTOR45}, or its central
\texttt{PURE5-DP-FULLYCROSSED-ENTROPYGAP45} restriction, is the first exact
residual.  Those objects remain valid parents.  The fully-crossed central child
has now been reduced further to a primitive B\'ezout-line correlation and then
to a canonical rank-one M\"obius/five-defect theorem.  No theorem in this
section closes that final correlation.
\end{status}

\subsection{From fully-crossed incidence to the double-GCD B\'ezout line}
In the entropy gap $H<P<U$, the common-$p$ and common-$d$ off-diagonal
strata are empty, so the source is fully crossed:
\[
 d_1\neq d_2,\qquad p_1\neq p_2.
\]
The large-prime divisor lemma can be sharpened on the clean cell: a structured
integer $N\ll Q/g$ has at most two prime divisors in the $P$-box, since three
such divisors would give $N\ge P^3>Q$.  We therefore bank
\[
 \boxed{\texttt{FULLCROSS-LARGEPRIME-DIVISOR-MULT2-45: PASS}}.
\]
In the subrange $g\gg Q/P^2$ the $P$-scale divisor is unique.

Write the common factor on the $d$ side as
\[
 d_1=ga,\qquad d_2=gb,\qquad (a,b)=1,
\]
and the defect-side common divisor as
\[
 v_1=sw_1,\qquad v_2=sw_2,\qquad (w_1,w_2)=1.
\]
The exact difference relations force $s\mid j$; with $j=sk$ the fully-crossed
system becomes
\begin{equation}
 \boxed{
 \begin{cases}
  g\ell A-usw_1=2,\\
  A_2-A_1=usk,\\
  w_2-w_1=g\ell k,
 \end{cases}}
 \label{eq:d16-doublegcd-system}
\end{equation}
where $A=ap$ denotes the M\"obius-prime factor.  Hence
\[
 \boxed{\texttt{FULLCROSS-DEFECT-GCD-ROUTER45: PASS}}.
\]
\[
 \boxed{\texttt{FULLCROSS-DOUBLEGCD-BEZOUT45: PASS}}.
\]

For fixed small $(g,s)$ choose a base solution
\[
 g\ell A_0-usw_0=2
\]
and parameterise the primitive line by
\[
 A_t=A_0+us\,t,\qquad
 w_t=w_0+g\ell t,\qquad
 T_{g,s}\asymp \frac{H}{gs}.
\]
Clean coprimality and the size inequalities imply that a fixed large prime
$p$ and a fixed M\"obius cofactor $a$ each occur at most once on a line:
\[
 \boxed{\texttt{DP-FACTOR-INJECTIVE-BEZOUTLINE45: PASS}}.
\]
The absolute family scale is
\[
 UR\,T_{g,s}^2\asymp \frac{QV}{g^2s^2}.
\]
Thus the tail $gs>(\log X)^{B_A}$ is summable with arbitrary logarithmic
saving:
\[
 \boxed{\texttt{DOUBLEGCD-LARGE-GS-TAIL45: PASS}}.
\]
The primitive $g=s=1$ child, however, remains at the full natural $QV$ scale.
Consequently
\[
 \boxed{\texttt{FULLCROSS-RIGIDITY-ONLY-CLOSURE45: FAILED ROUTE}}.
\]
This is a proof firewall: GCD reduction, support sparsity, bounded divisor
multiplicity and factor injectivity do not supply sign cancellation.

\subsection{Source-preserving MRK5--MOYK5 reverse audit}
A separate reverse audit returns to the scalar source before the destructive
Cauchy step.  The MRK5 Poisson skeleton is recovered in explicit-model form
and a micro-switch manufactures the source-faithful MOYK5 kernel
\[
 \boxed{S(k\overline B,-2h;m)},
 \qquad
 B\asymp Y^7,\quad m\asymp Q,\quad |h|,|k|\ll Q/Y.
\]
We bank
\[
 \boxed{\texttt{MRK5-POISSON-SKELETON: PASS}},
 \qquad
 \boxed{\texttt{MRK5-SOURCE-RECOVERED: PASS}},
\]
\[
 \boxed{\texttt{MOYK5-MICRO-SWITCH45: PASS}},
 \qquad
 \boxed{\texttt{MOYK5-SOURCE-KERNEL-MANUFACTURE45: PASS}}.
\]
The resulting theorem auditions are negative rather than closing:
\begin{itemize}
 \item the proposed Yang Case-5 landing has the wrong coefficient class;
 \item the inner Blomer--Pascadi dictionary can be written, but the physical
       lengths lie outside the range in which the global application pays;
 \item a further character/completion transform returns the QK5 signed
       representation rather than a smaller theorem.
\end{itemize}
Thus
\[
 \boxed{\texttt{MOYK5-QK5-REPRESENTATION-ANTILOOP45: PASS}}
\]
and MOYK5 does not supersede the B\'ezout-line residual.  This is consistent
with the published scope of Blomer--Pascadi's all-moduli bilinear Kloosterman
bound, whose strongest stated critical-range saving is at square-root length
\cite{BlomerPascadi2026}.

\subsection{The semiprime coefficient is literally M\"obius times a large-prime selector}
On a clean squarefree cell define
\[
 L_{D,P}(n)
 :=
 \sum_{\substack{p\mid n\\p\sim P,\ n/p\sim D}}\log p.
\]
Then
\[
 \beta_{D,P}(n)
 =\sum_{\substack{dp=n\\d\sim D,\ p\sim P}}\mu(d)\log p
 =-\mu(n)L_{D,P}(n),
\]
because $\mu(n/p)=-\mu(n)$ for squarefree $n$ and $p\mid n$.  The
large-prime multiplicity is $O(1)$, and in the central clean cell is at most
two.  We therefore bank
\[
 \boxed{\texttt{BETA-DP-TO-MU-LARGEPRIME45: PASS}}.
\]
The parity carrier in the primitive line is no longer an opaque semiprime
coefficient: it is literally a two-point M\"obius correlation.

\subsection{Canonical primitive rank-one form}
For the primitive $g=s=1$ child choose
\[
 \ell A_0-u w_0=2,\qquad (u,\ell)=1,
\]
and set
\[
 A_t=A_0+ut,\qquad w_t=w_0+\ell t,\qquad |t|\ll H.
\]
The exact residual is
\begin{equation}
\boxed{
\begin{aligned}
 \mathcal R_{\rm R1}
 :=
 \sum_{\substack{u\sim U,\ \ell\sim R\\(u,\ell)=1}}
 \sum_t\sum_{\substack{k\ne0\\|t+k|\ll H}}
 &\mu(A_t)\mu(A_{t+k})\\
 &\times L_{D,P}(A_t)L_{D,P}(A_{t+k})\\
 &\times b_5(w_t)\overline{b_5(w_{t+k})}\,
 \mathcal W_{u,\ell,t,k}.
\end{aligned}}
\label{eq:d16-rankone-residual}
\end{equation}
The required bound is
\begin{equation}
 \boxed{|\mathcal R_{\rm R1}|\ll_A QV(\log X)^{-A}.}
 \label{eq:d16-rankone-target}
\end{equation}
We name \eqref{eq:d16-rankone-target}
\[
 \boxed{\texttt{RANKONE-BEZOUT-AFFINE-ELLIOTT5D45}}.
\]
This is equivalent to the primitive small-$(g,s)$ child of
\texttt{PURE5-DP-DOUBLEGCD-BEZOUT-LINE45}; the new name records the exact
M\"obius and shift structure rather than evading the old residual.

There is also a useful geometric formulation.  Put
\[
 M_t=
 \begin{pmatrix}
 A_t&u\\
 w_t&\ell
 \end{pmatrix}.
\]
Then
\[
 \det M_t=2,\qquad
 M_t=M_0
 \begin{pmatrix}1&0\\t&1\end{pmatrix}.
\]
Thus the source runs along a discrete unipotent orbit in the determinant-$2$
integral matrices.  Line Fourier analysis, MOYK5 reciprocal kernels, the QK5
character parent and the post-Cauchy B\'ezout graph are different coordinates
on the same orbit; this explains the repeated representation anti-loops.

\subsection{Rank-one shift geometry and the failure of a direct all-shift Chowla splice}
For an off-diagonal displacement define
\[
 r=uk,\qquad s=\ell k.
\]
Since $(u,\ell)=1$,
\[
 k=(r,s),\qquad u=\frac r{(r,s)},\qquad
 \ell=\frac s{(r,s)}.
\]
Hence the map $(u,\ell,k)\mapsto(r,s)$ is injective on the clean positive
sector:
\[
 \boxed{\texttt{RANKONE-SHIFT-BIJECTION45: PASS}}.
\]
The actual shift family has size
\[
 URH=X,
\]
whereas the independent rectangle $r\ll Q$, $s\ll V$ has size
\[
 QV=XH.
\]
Since
\[
 H\ge X^{5/18},
\]
the physical shifts occupy only an $H^{-1}$-density rank-one subset of the
full rectangle.  Embedding this family into a complete-shift average would
therefore incur a fixed-power tax.

This blocks a direct splice from the published averaged-Chowla theorem of
Matom\"aki--Radziwi\l\l--Tao, which gains cancellation after averaging over a
complete family of shifts \cite{MatomakiRadziwillTao2015}.  The logarithmic
correlation structure theorem of Tao--Ter\"av\"ainen likewise concerns
logarithmically averaged fixed-shift correlations \cite{TaoTeravainen2017}.
The recent quantitative growing-shift result of Guo remains logarithmically
weighted and only treats polylogarithmic shifts \cite{Guo2026}.  Finally,
short-arithmetic-progression variance for multiplicative functions
\cite{KlurmanMangerelTeravainen2019} is not a vector-valued estimate for the
shift-dependent five-defect projector.  We therefore bank the source-specific
negative dictionary
\[
 \boxed{\texttt{ALLSHIFT-CHOWLA-TO-RANKONE45: FALSE AS A DIRECT SPLICE}}.
\]
This statement does \emph{not} prove that
\texttt{RANKONE-BEZOUT-AFFINE-ELLIOTT5D45} is equivalent to ordinary
fixed-shift Chowla.  The physical problem has extra $(u,\ell,k)$ averaging,
fixed determinant $2$, and a five-fold centred source, but that averaging lies
on a power-thin rank-one variety.

\subsection{Fresh fixed-denominator audition and death certificate}
The explicit MOYK5 reciprocal phase can be arranged as
\[
 e_{R_fn}(2h\overline m),
 \qquad
 R_fn\asymp Y^8,\quad m\asymp Q,\quad |h|\ll Q/Y.
\]
This superficially resembles Wright's trilinear Kloosterman-fraction theorem
with a partially fixed denominator factor \cite{Wright2026}.  A deliberately
generous audit over every admissible split
\[
 R_f=Y^\rho,\qquad n=Y^{8-\rho}
\]
gives a best total exponent $Y^{27/2+o(1)}=X^{3/2+o(1)}$, whereas the physical
scalar target is $Y^9=X$.  Hence even the most favourable split leaves the
fixed-power deficit $Y^{9/2}=X^{1/2}$:
\[
 \boxed{\texttt{WRIGHT-PARTFIX-MOYK5-45: POWER NONCLOSING}}.
\]
Together with the literal Blomer--Pascadi range audit, this strongly suggests
that further fixed-modulus Kloosterman theorem shopping is not the immediate
closure mechanism.

\subsection{Low pretentious profiles are eliminated}
Let $\chi$ be primitive with
\[
 \operatorname{cond}\chi\le(\log X)^B,\qquad |\tau|\le(\log X)^B.
\]
For one centred defect set
\[
 \mathcal D_i(\chi,\tau)
 =\sum_{r\sim Y}\delta_i(r)\chi(r)r^{i\tau}.
\]
Under the standard smooth packet conventions, small-conductor prime number
theorems and partial summation give, for every $C>0$,
\[
 \mathcal D_i(\chi,\tau)\ll_{B,C}Y(\log X)^{-C}.
\]
For the semiprime source,
\[
 \mathcal B_{D,P}(\chi,\tau)
 =
 \left(\sum_{d\sim D}\mu(d)\chi(d)d^{i\tau}W_D(d)\right)
 \left(\sum_{p\sim P}\log p\,\chi(p)p^{i\tau}W_P(p)\right).
\]
In the entropy gap $D>H\ge X^{5/18}$, the M\"obius factor has
arbitrary-log Fourier cancellation; for nonprincipal small-conductor $\chi$,
the prime factor is controlled by Siegel--Walfisz.  Thus
\[
 \mathcal B_{D,P}(\chi,\tau)\ll_{B,C}Q(\log X)^{-C}.
\]
We record this internal analytic child as
\[
 \boxed{\texttt{LOW-PRETENTIOUS-PROFILE45: PASS}}.
\]
It rules out a common low-conductor, low-frequency profile as the cause of a
large simultaneous line overlap.  It is not a full inverse theorem: the
surviving branch may be high-conductor, high-frequency, or distributed.

\subsection{The actual first unproved implication}
The nearest plausible architecture is now an inverse theorem:
\[
\boxed{
\begin{array}{c}
\text{large rank-one correlation}\\
\downarrow\\
\begin{cases}
\text{common low profile } \chi n^{it},&
 \operatorname{cond}\chi,|t|\le\log^B X,\\
\text{or a quantitatively controlled high-profile packet.}
\end{cases}
\end{array}}
\]
The first branch is eliminated above.  The unresolved implication is that a
large value of \eqref{eq:d16-rankone-residual} forces the second branch with
enough entropy or multiplicity penalty to sum to
$QV(\log X)^{-A}$.  This can be described as a rank-one entropy decrement or a
vector-valued inverse large sieve, but neither label denotes a currently
published theorem matching the physical source.

The best presently banked estimate has size
\[
 QV(\log X)^{C_0},
\]
against the target \eqref{eq:d16-rankone-target}.  Thus there is
\emph{no power deficit}; the remaining analytic deficit is arbitrary
logarithmic cancellation on the rank-one B\'ezout family.

\subsection{Controlling ledger and updated Gate--1B DAG}
The strict Draft~16 ledger is
\begin{center}
\small
\begin{tabular}{p{0.61\textwidth}p{0.27\textwidth}}
\toprule
Item & Status\\
\midrule
\texttt{PURE5-RAW-SOURCE45} & \textbf{PASS}\\
\texttt{DP-LINE-ENTROPY-TRICHOTOMY45} & \textbf{PASS}\\
\texttt{FULLCROSS-LARGEPRIME-DIVISOR-MULT2-45} & \textbf{PASS}\\
\texttt{FULLCROSS-DEFECT-GCD-ROUTER45} & \textbf{PASS}\\
\texttt{FULLCROSS-DOUBLEGCD-BEZOUT45} & \textbf{PASS}\\
\texttt{DP-FACTOR-INJECTIVE-BEZOUTLINE45} & \textbf{PASS}\\
\texttt{DOUBLEGCD-LARGE-GS-TAIL45} & \textbf{PASS}\\
\texttt{FULLCROSS-RIGIDITY-ONLY-CLOSURE45} & \textbf{FAILED ROUTE}\\
\texttt{MRK5-SOURCE-RECOVERED} & \textbf{PASS}\\
\texttt{MOYK5-SOURCE-KERNEL-MANUFACTURE45} & \textbf{PASS}\\
\texttt{MOYK5-QK5-REPRESENTATION-ANTILOOP45} & \textbf{PASS}\\
\texttt{BETA-DP-TO-MU-LARGEPRIME45} & \textbf{PASS}\\
\texttt{RANKONE-SHIFT-BIJECTION45} & \textbf{PASS}\\
\texttt{ALLSHIFT-CHOWLA-TO-RANKONE45} & \textbf{FALSE AS DIRECT SPLICE}\\
\texttt{WRIGHT-PARTFIX-MOYK5-45} & \textbf{POWER NONCLOSING}\\
\texttt{LOW-PRETENTIOUS-PROFILE45} & \textbf{PASS}\\
\texttt{RANKONE-BEZOUT-AFFINE-ELLIOTT5D45} & \textbf{FIRST EXACT ANALYTIC OPEN}\\
\texttt{PURE5-COMPARISON-MAINTERM-PIN} & \textbf{SOURCE OPEN}\\
\texttt{D2-RAPID-DECAY-PIN} & \textbf{PASS FOR FIXED SMOOTH PACKETS}\\
\texttt{PURE5-DP-SIGNED45} & \textbf{OPEN}\\
\texttt{NEARPRIM-DP-SIGNED45} & \textbf{OPEN}\\
$r>1$ square-character family & \textbf{OPEN}\\
CSTAR--CNW transition strip & \textbf{OPEN}\\
QK56 full parent / exhaustiveness & \textbf{OPEN}\\
\texttt{GATE1B} & \textbf{OPEN}\\
\bottomrule
\end{tabular}
\end{center}

The shortest honest dependency graph is now
\[
\boxed{
\begin{array}{c}
\text{PURE5 scalar source}\\
\downarrow\\
\text{entropy gap + double GCD + B\'ezout geometry}\quad\mathbf{PASS}\\
\downarrow\\
\text{large-}gs\text{ tail}\quad\mathbf{PASS}\\
\downarrow\\
\boxed{\texttt{RANKONE-BEZOUT-AFFINE-ELLIOTT5D45}}\\
\downarrow\\
\text{high-profile rank-one inverse theorem}\\
\downarrow\\
\texttt{PURE5-DP-SIGNED45}
+\text{ comparison/main-term pin}\\
\downarrow\\
\texttt{NEARPRIM-DP-SIGNED45}\\
\downarrow\\
r>1\text{ conductor family}+\text{transition strip}+\text{triangular recursion}\\
\downarrow\\
\texttt{QK56-FULL-PARENT}\\
\downarrow\\
\texttt{QK56-EXHAUSTIVENESS}\\
\downarrow\\
\text{finite source reassembly}\\
\downarrow\\
\texttt{GATE1B CLOSED}.
\end{array}}
\]
Every arrow after the first exact residual remains conditional on the stated
analytic/source antecedents.  In particular, this section does not prove
Gate~1B, Full Type II, or twin-prime infinitude.


\section{Working missing-literature and source-definition bank}
\label{app:missing-source-bank}

\begin{status}[Editorial working ledger]
This appendix is a working retrieval ledger rather than part of the proof of
any theorem above.  It is retained so that source and literature gaps can be
cleared systematically in a later final audit.  In particular, a missing
\emph{programme source definition} is not to be relabelled as a missing
published theorem.
\end{status}

\subsection{ML/SOURCE-1: Gate-1B physical-source definition gap}
\label{app:ml-source-1}

The current archive search recovers substantial finite compiler mathematics,
but does not recover the two physical source definitions needed to instantiate
the remaining Gate-1B compiler.  We therefore bank a single item
\[
   \boxed{\texttt{ML/SOURCE-1 --- Gate1B physical-source definition gap}}.
\]
This is principally a \emph{source-definition} gap, not a missing-paper gap.

\paragraph{Recovered finite/source algebra.}
The supplied archives contain the switched operator and its exact multiplier
reindexing, the $\mu*\Lambda$ divisor opening with ordered divisor-pair
multiplicity retained, and the higher-prime-power/repeated-prime/generic
switched partition.  They also contain the conditional Q5 support identity:
if the relevant coefficient is literally a Dirichlet convolution $c=a*b$,
then the generic switched term expands onto the determinant shell
\[
        mn+2=dpr.
\]
Likewise, the formal nine-coordinate ANOVA/distributive compiler and the
finite packet-reassembly compiler are available.  None of these finite
identities, however, identifies the actual physical nine-coordinate
coefficient.

\paragraph{Missing source definition 1: actual switched $R9$ cell.}
The first source-open item is
\[
  \boxed{\texttt{ACTUAL-SWITCHED-R9-CELL-DICTIONARY45}}.
\]
The formal interface permits a relation of the form
\[
  c_9(t)=\kappa_j\,(\alpha_j*\beta_{9-j})(t)+E_j(t),
\]
but the supplied source does not prove that the actual physical $c_9$ has
this form.  A source-safe dictionary must identify, for every actual packet,
\begin{itemize}
  \item the literal sequences $\alpha_j$ and $\beta_{9-j}$;
  \item labelled versus symmetrised source coordinates and their dyadic
        ranges;
  \item the M\"obius, prime and model scalars;
  \item ordered-factorisation multiplicity and repeated-prime correction;
  \item smooth/Perron weights, the original prefactor and physical target;
  \item exceptional and degeneracy routing.
\end{itemize}
Until this dictionary exists, the exact nine-coordinate convolution compiler
is only conditional on its source antecedent.

\paragraph{Missing source definition 2: shifted $TT^*$ physical source.}
The second source-open item is
\[
  \boxed{\texttt{SHIFT-TTSTAR-LITERAL-SOURCE-PIN}}.
\]
The archive contains four-cycle matrix, trace, determinant and discriminant
algebra, and counterguards against treating Cauchy copies as independent
physical multipliers.  It does not contain the literal pre-support-enlargement
$TT^*$ source formula from which the cyclic multipliers should be derived.
Before \texttt{SHIFT-SOURCE-LINKED-CHAR45} is a fully source-posed theorem,
the source pin must specify
\begin{itemize}
  \item the literal definitions and ranges of $B_1,\ldots,B_4$;
  \item the exact cyclic law $\Theta(B_i,B_{i+1})$;
  \item all $h_i$ ranges and source weights;
  \item the prefactor, normalisation and physical target;
  \item the $D=0$, non-unit, shared-prime, repeated-prime and collision strata;
  \item labelled tuple multiplicity and the exact equality/error across any
        support enlargement.
\end{itemize}

\paragraph{Why literature search cannot by itself close this item.}
Published determinant-character, trace-function and Kloosterman estimates can
be matched only \emph{after} the physical coefficient and multiplier family
are defined.  They cannot determine a private programme coefficient such as
$c_9$, the variables $B_i$, the cyclic map $\Theta(B_i,B_{i+1})$, or the
packet multiplicity.  Ford--Maynard's published sieve framework likewise does
not define these programme-specific objects \cite{FordMaynard2024}.
Consequently the next action for this bank item is source archaeology or a
literal source derivation, not broader theorem shopping.

\paragraph{Current status and later final-audit checklist.}
We retain
\[
\begin{array}{ll}
 \text{switched finite compiler} & \text{pass},\\
 \text{nine-coordinate convolution compiler} & \text{conditional pass},\\
 \text{finite Gate-1B reassembly compiler} & \text{pass},\\
 \text{actual switched $R9$ source dictionary} & \text{source open},\\
 \text{shifted $TT^*$ physical source} & \text{source open},\\
 \text{Gate 1B} & \text{open}.
\end{array}
\]
When the final missing-literature/source audit is compiled, this item should
be cleared only by one of the following: (i) an authoritative source file
containing the missing literal definitions; (ii) a derivation from an earlier
physical source which proves them exactly, including multiplicity and
normalisation; or (iii) a proof that the corresponding packet is absent and
therefore need not be routed.  A published black box matching only the later
schematic kernel is not sufficient.

\printbibliography
\end{document}
